\documentclass{article}
\usepackage{graphicx}
\usepackage[utf8]{inputenc}
\usepackage[T1]{fontenc}
\usepackage[english, catalan, spanish]{babel} 
\usepackage{geometry}
\usepackage{fancyhdr}
\usepackage{amsmath, amssymb, amsfonts}
\usepackage{listings}
\usepackage{xcolor}
\usepackage{hyperref}
\usepackage{float}
\usepackage{booktabs}
\usepackage{enumitem}
\usepackage{titlesec}
\usepackage{pgfplots}
\usepackage[absolute]{textpos}
\pgfplotsset{compat=1.17}
\usepackage{ragged2e}
\usepackage{setspace}
\usepackage{rotating}
\usepackage{pdflscape}
\usepackage{eurosym}
\usepackage{amsmath}
\usepackage[backend=biber,style=numeric,sorting=none]{biblatex}
\addbibresource{Mi_biblioteca.bib}

\lstdefinelanguage{terraform}{
    keywords={resource,data,provider,variable,output,locals,module,terraform,required_providers},
    keywordstyle=\bfseries,
    sensitive=true,
    comment=[l]{\#},
    comment=[l]{//},
    comment=[s]{/*}{*/},
    string=[b]",
    morestring=[b]',
}

\pagestyle{fancy}
\fancyhf{}
\fancyhead[L]{Informe de seguimento}
\fancyfoot[C]{\thepage}
\renewcommand{\headrulewidth}{0.4pt}

\title{\textbf{INFORME DE SEGUIMIENTO DEL PROYECTO}}
\author{Omar Antonio Cornejo Vargas}
\date{March 2026}

\begin{document}

    \begin{titlepage}
    \centering

    \begin{minipage}{0.1\textwidth}
        \centering
        \includegraphics[width=\textwidth]{logo-upc.png}
    \end{minipage}
    \hfill
    \begin{minipage}{0.88\textwidth}
        \centering
        \includegraphics[width=\textwidth]{logo-fiblletres-upc-color.png}
    \end{minipage}

    \vspace{1cm}

    {\Large \textbf{FACULTAD DE INFORMÁTICA DE BARCELONA (FIB)} \par}
    \vspace{0.2cm}
    {\large UNIVERSIDAD POLITÉCNICA DE CATALUÑA (UPC) \par}
    \vspace{0.2cm}
    {\large ESPECIALIDAD EN TECNOLOGÍAS DE LA INFORMACIÓN \par}

    \vspace{1.5cm}

    {\huge \textbf{DISEÑO DE UN SISTEMA VISUAL PARA LA CREACIÓN Y DESPLIEGUE DE ARQUITECTURAS MEDIANTE IAC} \par}

    \vspace{1.5cm}

    {\Large \textbf{TRABAJO FINAL DE GRADO} \par}
    \vspace{0.3cm}
    {\large INFORME DE SEGUIMIENTO DEL PROYECTO \par}

    \vspace{2cm}

    {\LARGE \textbf{Omar Antonio Cornejo Vargas} \par}

    \vfill
    
    {\large \textbf{Director:} Marc Catrisse i Pérez \par}
    {\large \textbf{Tutora GEP:} Olga Pons Peregort \par}
    {\large \textbf{Titulación:} Ingeniería Informática \par}
    
    \vspace{1cm}
    
    {\large 21 de mayo de 2026 \par}

    \end{titlepage}



    \newpage
    \renewcommand{\contentsname}{\Huge \textbf{ÍNDICE}}
    {\large
    \tableofcontents
    }

    \newpage
    \section{Introducción y contextualización}
    {                     
        \setstretch{1.2}    
        \setlength{\parskip}{1em}

        \justify 
        Este proyecto se desarrolla en el marco de la \textbf{Facultad de Informática de Barcelona (FIB)} de la \textbf{Universidad Politécnica de Cataluña (UPC)}, en el contexto de la especialización en Tecnologías de la Información. La propuesta se centra en aplicar los conocimientos adquiridos durante la titulación y aquellos aprendidos de forma autónoma, con el objetivo de abordar un problema real en el ámbito del \textbf{Cloud Computing} \cite{noauthor_que_2024} y la \textbf{Infraestructura como Código (IaC)} \cite{noauthor_que_nodate-4}.

        \subsection{Contextualización}
        \justify
        
        El \textbf{Cloud Computing}, o computación en la nube permite a los usuarios acceder a los servicios de manera flexible, escalable y bajo demanda, sin la necesidad de mantener infraestructura local extensa. Esta modalidad simplifica la gestión de sistemas complejos, optimiza recursos y facilita la implementación rápida de soluciones en entornos dinámicos y cambiantes.

        \justify
        El Cloud Computing se ha convertido en un pilar fundamental para empresas de todos los tamaños, permitiendo lanzar productos y servicios con mayor rapidez y eficiencia, reducir costes de infraestructura y adaptarse a la demanda cambiante del mercado. Además, facilita la colaboración  remota, la integración de tecnologías emergentes como inteligencia artificial, análisis de datos y aplicaciones distribuidas. Su adopción creciente transforma no solo la manera en que se gestionan los sistemas de información, sino también la forma de organizar los procesos empresariales y tecnológicos en un mundo cada vez más digitalizado.

        \subsection{Términos clave}
        \justify
        A continuación se presentan los conceptos fundamentales que se utilizarán a lo largo de este proyecto. Comprender estos términos es esencial para entender tanto el contexto como la solución propuesta.
        
        \begin{itemize}
            \item \textbf{Cloud Computing (Computación en la nube):} El cloud computing es el acceso bajo demanda a recursos informáticos (servidores físicos o virtuales, almacenamiento de datos, capacidades de red, herramientas de desarrollo de aplicaciones, software, plataformas analíticas con IA y más) a través de Internet con precios de pago por uso.
            
            \item \textbf{Infraestructura como Código (IaC):} La infraestructura como código en DevOps es una herramienta poderosa para agilizar la configuración de la infraestructura, especialmente en entornos de nube. IaC automatiza el proceso de diseño y creación de la infraestructura.
            
            \item \textbf{Proveedor hiperescalar:} Los hyperscalers ofrecen servicios de cloud computing y de gestión de datos para las empresas que requieren una amplia infraestructura para el procesamiento y el almacenamiento de datos a gran escala. Si bien no hay un estándar universal para definir a los hyperscalers, los distribuidores de nube más importantes, como Amazon Web Services \cite{noauthor_que_nodate-6}, Google Cloud \cite{noauthor_descripcion_nodate},Microsoft Azure \cite{noauthor_que_nodate-5}, IBM Cloud \cite{noauthor_ibm_nodate} y Alibaba Cloud \cite{noauthor_alibaba_nodate}, se ajustan a la descripción \cite{noauthor_hyperscalers_nodate} 
            
            \item \textbf{Escalabilidad:} La escalabilidad es la capacidad de un sistema, proceso o recurso para adaptarse y manejar un aumento en las demandas del usuario sin perder su rendimiento o funcionalidad \cite{noauthor_escalabilidad_nodate}.
            
            \item \textbf{Portabilidad:} La portabilidad es una característica del software, siendo una forma de reutilización que ayuda a evitar el “bloqueo” de ciertos entornos operativos, p.ej. proveedores de nube, sistemas operativos o vendedores \cite{noauthor_portabilidad_nodate}.

            \item \textbf{Sostenibilidad:} La sostenibilidad informática va más allá de adoptar tecnologías ecológicas. Implica repensar cómo se gestionan los recursos digitales desde el hardware hasta la nube para optimizar el consumo, extender la vida útil de los equipos y fomentar la innovación sostenible en IT \cite{noauthor_que_nodate-3}.
            
            \item \textbf{Open Source:} El software open source es un código diseñado de manera que sea accesible al público: todos pueden ver, modificar y distribuir el código de la forma que consideren conveniente \cite{noauthor_que_nodate-2}.

            
        \end{itemize}
        
        \justify
        Estos términos serán los más recurrentes en el desarrollo de la propuesta y son la base para entender los desafíos que enfrenta la industria y la relevancia de la solución planteada en este proyecto.

        \subsection{Problema}
        \justify
        A pesar de sus ventajas, los equipos que trabajan en entornos Cloud enfrentan barreras importantes para diseñar y desplegar infraestructuras de manera eficiente. Las herramientas existentes suelen ser fragmentadas, presentan curvas de aprendizaje elevadas y limitaciones de interoperabilidad, lo que complica los procesos y reduce la productividad. Además, los distintos proveedores hiperescalares, como \textbf{Amazon Web Services (AWS)}, \textbf{Google Cloud Platform (GCP)} y \textbf{Microsoft Azure}, ofrecen servicios con características, nomenclaturas y modelos de facturación diferentes, lo que dificulta la estandarización de procesos y la portabilidad de aplicaciones entre plataformas.
        
        \justify
        La gestión de entornos en la nube requiere priorizar la escalabilidad y la sostenibilidad a largo plazo, garantizando que los servicios puedan evolucionar y adaptarse a nuevas demandas. Esto hace necesario que los proyectos incorporen desde el inicio herramientas que proporcionen flexibilidad y adaptabilidad en la infraestructura, permitiendo un crecimiento orgánico y eficiente de los sistemas.
        
        \justify
        Ante este escenario, se hace evidente la necesidad de soluciones más accesibles, intuitivas y de código abierto, que combinen facilidad de uso con agilidad en la implementación, contribuyendo a reducir la complejidad y mejorar la eficiencia en el diseño y despliegue de infraestructuras en la nube.
        \par 

        
        \subsection{Actores implicados}
        \label{sec:actores_implicados}
        \justify
        El producto está dirigido y utilizado por distintos actores dentro del ecosistema Cloud, y sus resultados impactan a varios niveles:
        
        \begin{itemize}
            \item \textbf{Proveedores de servicios Cloud:} Empresas como \textbf{Amazon Web Services (AWS)}, \textbf{Google Cloud Platform (GCP)} y \textbf{Microsoft Azure}. Son los encargados de ofrecer los recursos necesarios para desplegar y operar las infraestructuras. Se benefician indirectamente al facilitar el uso eficiente de sus servicios y fomentar la adopción de la nube.
        
            \item \textbf{Usuarios y organizaciones:} Usuarios finales, equipos de desarrollo, operaciones y administración de sistemas que utilizarán el producto para diseñar, implementar y mantener aplicaciones e infraestructuras. Son los usuarios principales y se benefician directamente de la automatización, estandarización y eficiencia que el sistema proporciona.
        
            \item \textbf{Comunidad y ecosistema Open Source:} Contribuye al desarrollo de herramientas, buenas prácticas y documentación, facilitando la colaboración y el acceso a soluciones abiertas. Se beneficia del intercambio de conocimiento y de la mejora continua de las soluciones que el producto integra o promueve.
        \end{itemize}
        
       
    }

    \newpage
    \section{Relevancia y Justificación}

    {                    
        \setstretch{1.2}            
        \setlength{\parskip}{1em}   

        \justify 
        El diseño y despliegue de infraestructuras en la nube cuenta actualmente con diversas herramientas consolidadas en la industria. Entre ellas destacan soluciones de Infraestructura como Código (IaC) como Terraform \cite{noauthor_que_2025}, que permiten definir y gestionar recursos mediante archivos de configuración, así como plataformas de modelado visual como Brainboard \cite{noauthor_brainboard_nodate}. Asimismo, los principales proveedores hiperescaladores, como Amazon Web Services (AWS), Google Cloud Platform (GCP) y Microsoft Azure, ofrecen herramientas propias que permiten visualizar y gestionar los recursos desplegados dentro de sus respectivos entornos.
        \par 
    
        \justify 
        A pesar de la existencia de estas soluciones, cada una presenta limitaciones que dificultan cubrir simultáneamente criterios de accesibilidad, independencia de proveedor y facilidad de diseño. Las herramientas basadas en código, como Terraform, proporcionan un alto grado de control y automatización, pero requieren conocimientos técnicos avanzados y familiaridad con lenguajes de configuración específicos, lo que puede suponer una barrera de entrada para perfiles junior, estudiantes o usuarios que se inician en el ámbito del Cloud Computing. Por otro lado, las plataformas que ofrecen modelado visual de infraestructuras suelen estar asociadas a modelos de licencia propietarios o a costes de uso, lo que limita su accesibilidad y dificulta su utilización en entornos educativos o de experimentación.
        \par 
    
        \justify 
        Asimismo, las herramientas proporcionadas por los propios proveedores cloud permiten visualizar y administrar infraestructuras ya desplegadas, pero no están diseñadas principalmente para facilitar el diseño conceptual de arquitecturas desde cero ni para trabajar desde una perspectiva independiente del proveedor. Este enfoque puede favorecer situaciones de dependencia con proveedores y reducir la portabilidad de las arquitecturas diseñadas entre distintas plataformas cloud.
        \par 

        \justify
        Desde una perspectiva educativa y económica, este enfoque también aporta un valor añadido relevante. Una herramienta open source y orientada al aprendizaje facilita la experimentación con arquitecturas cloud sin necesidad de asumir costes asociados a licencias comerciales ni depender de plataformas propietarias. Esto favorece su utilización en contextos académicos, en procesos de formación y en entornos de prototipado, contribuyendo a reducir las barreras de acceso al aprendizaje del Cloud Computing.
        \par
        
        \justify
        En este contexto, la simple adaptación de una herramienta existente no resulta suficiente para cubrir de forma conjunta estos requisitos. El desarrollo de una nueva solución permite plantear una herramienta que combine las ventajas de la Infraestructura como Código con un entorno visual intuitivo y accesible, manteniendo además un enfoque abierto y neutral respecto a los distintos proveedores cloud y orientado tanto al aprendizaje como a la experimentación tecnológica.
        \par

        \subsection{Comparativa con soluciones existentes}
        \justify
        Para situar la propuesta frente al estado del arte, la Tabla~\ref{tab:comparativa_soluciones}
        compara las soluciones más representativas en cinco ejes que reflejan los requisitos del
        proyecto: la disponibilidad de un \textbf{modelado visual}, la \textbf{independencia
        respecto al proveedor cloud}, la condición de \textbf{open source}, la capacidad de
        \textbf{generar IaC editable} por el usuario, y la posibilidad de \textbf{ejecutar
        íntegramente en local} sin depender de una plataforma remota.

        \begin{table}[H]
        \centering
        \small
        \begin{tabular}{lccccc}
        \toprule
        \textbf{Solución} & \textbf{Visual} & \textbf{Indep.\ proveedor} & \textbf{Open source} & \textbf{IaC editable} & \textbf{Local} \\
        \midrule
        Terraform CLI               & No      & Sí      & Sí      & Sí      & Sí \\
        Pulumi                      & No      & Sí      & Sí      & Sí      & Sí \\
        AWS CDK                     & No      & No      & Sí      & Sí      & Sí \\
        Brainboard                  & Sí      & Sí      & No      & Parcial & No \\
        Consolas AWS / GCP / Azure  & Parcial & No      & No      & No      & No \\
        \midrule
        \textbf{Solución propuesta} & \textbf{Sí} & \textbf{Sí} & \textbf{Sí} & \textbf{Sí} & \textbf{Sí} \\
        \bottomrule
        \end{tabular}
        \caption{Comparativa de soluciones existentes frente a la propuesta}
        \label{tab:comparativa_soluciones}
        \footnotesize Fuente: elaboración propia.
        \end{table}

        \justify
        La tabla evidencia que cada solución cubre un subconjunto de los criterios, pero
        ninguna los reúne simultáneamente. Las herramientas basadas en código (Terraform,
        Pulumi~\cite{noauthor_infrastructure_nodate}, AWS CDK~\cite{noauthor_cloud_nodate})
        son potentes y locales, pero carecen de modelado visual y, en el caso
        de CDK, están ligadas a un proveedor concreto. Brainboard aporta el componente
        visual y la independencia de proveedor, pero es propietaria y depende de su
        plataforma remota. Las consolas oficiales de los hiperescaladores son visuales y
        gráficas, pero atan al usuario al proveedor y no generan IaC editable que pueda
        versionarse. La solución propuesta se sitúa precisamente en el hueco que estas
        herramientas dejan sin cubrir: un editor visual, open source, neutral respecto al
        proveedor, que produce IaC editable y se ejecuta en local.
        \par
    }


    \newpage
    \section{Alcance del proyecto}
    {                 
        \setstretch{1.2}            
        \setlength{\parskip}{1em}   
        \label{sec:Alcance}
        \subsection{Objetivo general}
        \justify
        Diseñar y desarrollar un sistema visual que permita modelar arquitecturas Cloud de 
        forma intuitiva y generar automáticamente la correspondiente definición de 
        Infraestructura como Código (IaC), promoviendo la accesibilidad, la interoperabilidad 
        entre proveedores y el enfoque open source.

        \justify
        No obstante, debido a las limitaciones temporales y al amplio alcance potencial del 
        proyecto, se plantea como un producto mínimo viable (MVP) \cite{noauthor_que_nodate-7}. 
        En consecuencia, se implementará un conjunto reducido y bien delimitado de 
        funcionalidades esenciales, priorizando aquellas que demuestren la viabilidad del 
        núcleo del sistema: la representación visual de recursos y la generación automática 
        de IaC.
        \par

        \subsection{Subobjetivos}
        \label{sec:Alcance_subobjetivos}
        \justify
        Una vez definido el objetivo general del proyecto, se establecen los siguientes 
        subobjetivos específicos, orientados a concretar y estructurar el desarrollo de 
        forma progresiva:

        \begin{itemize}
        \item Diseñar una interfaz gráfica que permita crear arquitecturas Cloud mediante 
        componentes visuales reutilizables e intuitivos, con soporte para la conexión 
        entre recursos.
        \item Implementar un mecanismo de traducción automática del diseño visual a código 
        de IaC estructurado y coherente, garantizando la correspondencia semántica entre 
        ambas representaciones.
        \item Definir una arquitectura de software modular y extensible que facilite la 
        incorporación futura de nuevos proveedores Cloud y tipos de recursos.
        \item Validar el sistema mediante casos de uso representativos que demuestren su 
        viabilidad técnica y su utilidad práctica en escenarios reales de despliegue.
        \end{itemize} 

        \subsection{Requisitos funcionales y no funcionales}
        \label{sec:Alcance_requisitos}
        \justify
        Desde el punto de vista funcional, el sistema debe permitir la creación, edición y 
        eliminación de recursos Cloud desde la interfaz gráfica, representar visualmente las 
        relaciones y dependencias entre dichos recursos, y generar automáticamente archivos 
        de configuración compatibles con herramientas de IaC. Asimismo, debe ofrecer 
        retroalimentación inmediata al usuario sobre la validez del diseño construido.

        \justify
        En cuanto a los requisitos no funcionales, la \textbf{usabilidad} y la 
        \textbf{accesibilidad} son aspectos centrales del diseño: la interfaz debe ser clara 
        e intuitiva, orientada tanto a perfiles técnicos como formativos, e inclusiva para 
        distintos tipos de usuario. La \textbf{portabilidad} exige un diseño agnóstico respecto 
        al proveedor Cloud, evitando dependencias que limiten su adopción. Por su parte, la 
        \textbf{escalabilidad} y la \textbf{mantenibilidad} condicionan las decisiones 
        arquitectónicas: el sistema debe estar preparado para incorporar nuevos servicios y 
        su código ha de ser estructurado y documentado. Todo ello bajo un modelo de 
        \textbf{licenciamiento open source} que garantice la libre distribución y 
        reutilización del proyecto.
        \par

        \subsection{Obstáculos y riesgos}
        \label{sec:Alcance_obstaculos}
        \justify
        El desarrollo del proyecto presenta varios riesgos que conviene identificar y acotar 
        desde el inicio. El principal reto técnico reside en la traducción precisa entre el 
        modelo visual y el código IaC: las diferencias estructurales entre proveedores Cloud 
        dificultan la definición de una abstracción común y coherente, y cualquier inconsistencia 
        en esta capa puede comprometer la corrección del código generado. En segundo lugar, la 
        complejidad inherente a la representación de dependencias entre recursos puede derivar 
        en un modelo de datos excesivamente complejo si no se delimita con rigor desde el diseño. 
        Finalmente, las limitaciones de tiempo y recursos propias de un Trabajo Final de Grado 
        condicionan el alcance de lo que es posible abordar.

        \justify
        Para mitigar estos riesgos, se adoptarán las siguientes estrategias: definir una capa de 
        abstracción común antes de implementar generadores específicos por proveedor; acotar el 
        MVP a un subconjunto representativo pero manejable de recursos y relaciones; y establecer 
        hitos parciales de revisión que permitan ajustar el alcance de forma controlada a lo largo 
        del desarrollo.
        \par
    }

    
    \newpage
    \section{Metodología de trabajo}
    {                
        \setstretch{1.2}            
        \setlength{\parskip}{1em}

        \justify
        Para el desarrollo de este proyecto se ha decidido combinar metodologías ágiles 
        adaptadas al contexto individual de un TFG. Se aplicará Scrum 
        \cite{kniberg_kanban_2010} para organizar sprints pactados con el director del
        proyecto, y Kanban para gestionar internamente las
        tareas de desarrollo de manera visual y flexible. Esta combinación responde a las 
        particularidades del proyecto: Scrum aporta la cadencia y los puntos de control 
        necesarios para un desarrollo supervisado, mientras que Kanban permite gestionar 
        la carga de trabajo diaria sin la rigidez de un equipo multidisciplinar, dado que 
        el proyecto lo desarrolla una sola persona.

        \subsection{Estructura del trabajo}
        \justify
        El proyecto se organiza en iteraciones o sprints de aproximadamente dos semanas, 
        cada uno con objetivos claros y entregables definidos. Al inicio de cada sprint 
        se planifican las tareas junto con el director, estableciendo metas alcanzables 
        y priorizando el trabajo más relevante según el estado del desarrollo. Al 
        finalizar cada iteración, se presenta el progreso obtenido y se recoge la 
        retroalimentación del director, lo que permite ajustar el alcance y reorientar 
        el siguiente sprint si fuera necesario.

        \justify
        Esta adaptación de Scrum difiere del marco original en varios aspectos deliberados: 
        aunque el proyecto lo desarrolla el autor solo, se distinguen tres perfiles 
        funcionales que se asume en distintos momentos del desarrollo: el de 
        \textbf{jefe de proyecto} (planificación y seguimiento), el de \textbf{ingeniero 
        de software} (diseño e implementación) y el de \textbf{redactor técnico} 
        (documentación y defensa). Esta distinción, en lugar de los roles formales de 
        Scrum (Product Owner, Scrum Master, desarrollador), permite asignar 
        responsabilidades de forma explícita y trazable a cada tarea. Las ceremonias 
        se simplifican a una reunión de revisión por sprint, y el backlog lo gestiona 
        directamente el propio autor. Estas adaptaciones son coherentes con las 
        limitaciones de un TFG individual y permiten mantener los beneficios del marco 
        ágil — iteratividad, visibilidad del progreso y capacidad de
        replanificación — sin incurrir en una sobrecarga de proceso innecesaria.
        \par

        \subsection{Herramientas de seguimiento y validación}
        \justify
        Las siguientes herramientas se utilizan no solo para organizar el trabajo, sino 
        específicamente para validar que se alcanzan los objetivos definidos en el alcance 
        del proyecto:

        \begin{itemize}
        \item \textbf{Control de versiones (GitHub) \cite{noauthor_acerca_nodate}:} 
        permite mantener trazabilidad de todos los cambios del código y verificar que 
        cada subobjetivo técnico se traduce en commits y ramas concretas. El historial 
        de commits actúa como evidencia objetiva del progreso real del proyecto.
        \item \textbf{Tablero Kanban (Trello) \cite{noauthor_que_nodate-8}:} organiza 
        las tareas en columnas de estado (pendiente, en progreso, completada), 
        permitiendo visualizar en todo momento si el avance del sprint es coherente 
        con los objetivos planificados y detectar bloqueos de forma temprana.
        \item \textbf{Cronograma y diagramas de Gantt:} establecen los hitos temporales 
        del proyecto y permiten comparar periódicamente el progreso real frente al 
        planificado, facilitando la detección de desviaciones y la toma de decisiones 
        sobre el alcance \cite{noauthor_que_nodate-9}.
        \end{itemize}

        \justify
        La combinación de estas herramientas crea un sistema de seguimiento en dos 
        niveles: el nivel de tarea (Trello) y el nivel de hito (Gantt y revisiones de 
        sprint), garantizando que tanto el día a día del desarrollo como los objetivos 
        globales del proyecto estén siempre bajo control.
        \par
    }

    \newpage
\section{Planificación temporal}
{                
    \setstretch{1.2}            
    \setlength{\parskip}{1em}
    
    \justify
    El proyecto se desarrolla entre el 9 de febrero y el 15 de junio de 2026, lo que
    representa aproximadamente 18 semanas de trabajo efectivo. La fase de gestión de
    proyecto ocupa las primeras semanas del calendario. Una vez concluida, las horas 
    restantes se destinan al desarrollo técnico del prototipo, las iteraciones de 
    revisión con el director y la redacción de la memoria y preparación de la defensa oral.

    \justify
    La planificación se organiza en sprints de dos semanas siguiendo la metodología ágil.
    La documentación se genera de forma continua en cada sprint, y las reuniones con el 
    director se celebran cada dos semanas durante toda la duración del TFG. 
    A continuación, se detallan los recursos necesarios para llevar a cabo correctamente 
    el proyecto.
    \par

    \subsection{Recursos del proyecto}

    \subsubsection{Recursos humanos}
        \justify
        Se identifican tres roles que se asumen en distintos momentos del proyecto:

        \begin{itemize}[leftmargin=*, itemsep=1pt]
        \item \textbf{Jefe de proyecto (JP).} Responsable de la planificación, seguimiento del
            cronograma y gestión de riesgos.
        \item \textbf{Ingeniero de software (IS).} Encargado del diseño arquitectónico, la
            implementación del prototipo y la investigación tecnológica.
        \item \textbf{Redactor técnico (RT).} Responsable de la redacción de la memoria, la
            documentación técnica y la preparación de la defensa oral.
        \end{itemize}

        \justify
        Adicionalmente, el director Marc Catrisse i Pérez actúa como colaborador externo,
        participando en las reuniones de seguimiento y validando las decisiones de diseño 
        y el alcance del prototipo.

        \subsubsection{Recursos materiales}

        \textbf{Hardware}

        \begin{itemize}[leftmargin=*, itemsep=1pt]
        \item Torre PC (AMD Ryzen 7 5800X, NVIDIA RTX 3060 12 GB, 32 GB RAM DDR4, 1 TB NVMe SSD)
        \item Monitor HP Omen
        \item Monitor Dell
        \item Teclado y ratón
        \end{itemize}

        \textbf{Software}

        \begin{itemize}[leftmargin=*, itemsep=1pt]
        \item \textbf{VS Code} — entorno de desarrollo integrado \cite{noauthor_visual_nodate}.
        \item \textbf{Rust + Tauri} — lenguaje e infraestructura del backend \cite{noauthor_rust_nodate} \cite{noauthor_que_2025-1}.
        \item \textbf{Vite + React, TypeScript} — stack del frontend \cite{noauthor_getting_nodate} \cite{noauthor_react_nodate} \cite{noauthor_javascript_nodate}.
        \item \textbf{GitHub} — control de versiones y gestión del repositorio
        \item \textbf{Trello} — tablero Kanban para el seguimiento de tareas
        \item \textbf{\LaTeX{}} — redacción de la memoria del proyecto \cite{noauthor_latex_nodate}.
        \item \textbf{Google Meet y correo electrónico} — comunicación con el director \cite{noauthor_google_nodate}.
        \end{itemize}

        \subsubsection{Recursos de infraestructura}

        \begin{itemize}[leftmargin=*, itemsep=1pt]
        \item \textbf{Espacio de trabajo:} domicilio del autor, utilizado como oficina 
        durante toda la duración del proyecto.
        \item \textbf{Electricidad:} consumo eléctrico de los dispositivos del puesto 
        de trabajo durante las 540 horas del proyecto.
        \item \textbf{Material de oficina:} suministros generales de papelería y 
        consumibles necesarios durante el desarrollo.
        \end{itemize}
    \par

    \subsection{Descripción de las tareas}
        \setstretch{1.2}            
        \setlength{\parskip}{1em}
    
        \justify
        Por cada tarea se indica descripción, estimación en horas, dependencias, recursos 
        materiales y rol humano responsable:
    
        \begin{itemize}[leftmargin=2em, itemsep=1pt, topsep=1pt]
        \item[R1] Ordenador personal con Linux e Internet
        \item[R2] \LaTeX{}
        \item[R3] Google Meet y correo electrónico
        \item[R4] Trello
        \item[R5] GitHub
        \end{itemize}

        \begin{itemize}[leftmargin=2em, itemsep=1pt, topsep=1pt]
        \item[JP] Jefe de proyecto
        \item[IS] Ingeniero de software
        \item[RT] Redactor técnico
        \end{itemize}


        \subsubsection{GEP – Gestión del Proyecto}
        \justify
        GEP (18 febrero 2026 → 18 marzo 2026), 90 h. 
        Elaboración del documento de gestión del proyecto, definiendo alcance y objetivos, 
        planificando tareas y tiempos, estimando costes y sostenibilidad, e integrando el 
        feedback de los tutores. \\
        \textit{Dep.: -- | Rec.: R1, R2, R3 | Rol: JP, RT}


        \subsubsection{DT – Desarrollo Técnico}
        \textit{Dedicación estimada: \textbf{450 h}. Comprende la investigación, el
        aprendizaje tecnológico y la implementación iterativa del prototipo.}
        
        \begin{itemize}[leftmargin=*, itemsep=2pt]
        
                \item \textbf{Sprint 1 – Investigación y definición inicial} 
        (18 marzo 2026 → 1 abril 2026) 60 h
        
            \begin{itemize}[leftmargin=1.2em, itemsep=2pt]
        
                \item \textbf{DT.1 – Investigación y estado del arte} 
                (15 h | Dep.: -- | Rec.: R1 | Rol: IS) \\
                Análisis de herramientas IaC (Terraform, Pulumi, CDK) y modelado visual 
                para detectar limitaciones y oportunidades.
        
                \item \textbf{DT.2 – Definición de requisitos y user stories} 
                (25 h | Dep.: DT.1 | Rec.: R1, R3 | Rol: JP, IS) \\
                Definición y refinamiento de requisitos e historias de usuario con 
                feedback del director.
        
                \item \textbf{DT.3 – Selección tecnológica} 
                (10 h | Dep.: DT.1 | Rec.: R1 | Rol: IS) \\
                Evaluación y justificación del stack tecnológico y lenguaje de 
                programación.
        
                \item \textbf{DT.4 – Seguimiento y documentación del sprint} 
                (10 h | Dep.: -- | Rec.: R1, R2, R3, R4 | Rol: JP, RT) \\
                Reuniones periódicas con el director, junto con el registro de 
                actividades, resultados y trazabilidad del sprint.
        
            \end{itemize}
        
            \item \textbf{Sprint 2 – Aprendizaje tecnológico y diseño arquitectónico} 
            (1 abril 2026 → 15 abril 2026) 65 h
            
            \begin{itemize}[leftmargin=1.2em, itemsep=2pt]
            
                \item \textbf{DT.5 – Aprendizaje de Rust, Tauri y React} 
                (30 h | Dep.: DT.3 | Rec.: R1, R5 | Rol: IS) \\
                Estudio de Rust, Tauri y React, con ejercicios y prototipos.
            
                \item \textbf{DT.6 – Diseño de la arquitectura del sistema} 
                (15 h | Dep.: DT.5 | Rec.: R1, R5 | Rol: IS) \\
                Definición de arquitectura modular (frontend React, lógica de grafos 
                y generación IaC en Rust).
            
                \item \textbf{DT.7 – Configuración del entorno} 
                (10 h | Dep.: DT.3 | Rec.: R1, R5 | Rol: IS) \\
                Preparación del entorno y estructura base del proyecto.
            
                \item \textbf{DT.8 – Seguimiento y documentación del sprint} 
                (10 h | Dep.: -- | Rec.: R1, R2, R3, R4 | Rol: JP, RT) \\
                Reuniones periódicas con el director, junto con el registro de 
                actividades, resultados y trazabilidad del sprint.
            
            \end{itemize}            

            \item \textbf{Sprint 3 – Primeras pruebas técnicas y mockups} 
            (15 abril 2026 → 29 abril 2026) 70 h
            
                \begin{itemize}[leftmargin=1.2em, itemsep=2pt]
            
                    \item \textbf{DT.9 – Canvas básico} 
                    (45 h | Dep.: DT.6, DT.7 | Rec.: R1, R4, R5 | Rol: IS) \\
                    Implementación inicial del editor visual con renderizado y 
                    drag-and-drop básico.
            
                    \item \textbf{DT.10 – Pruebas de traducción a IaC} 
                    (15 h | Dep.: DT.9 | Rec.: R1, R4, R5 | Rol: IS) \\
                    Primeras pruebas de generación de código HCL \cite{noauthor_hashicorphcl_2026} desde elementos 
                    visuales.
            
                    \item \textbf{DT.11 – Seguimiento y documentación del sprint} 
                    (10 h | Dep.: -- | Rec.: R1, R2, R3, R4 | Rol: JP, RT) \\
                    Reuniones periódicas con el director, junto con el registro de 
                    actividades, resultados y trazabilidad del sprint.
            
                \end{itemize}


            \item \textbf{Sprint 4 – Continuación de la implementación} 
            (29 abril 2026 → 13 mayo 2026) 90 h
            
            \begin{itemize}[leftmargin=1.2em, itemsep=2pt]
            
                \item \textbf{DT.12 – Editor visual completo} 
                (30 h | Dep.: DT.9 | Rec.: R1, R4, R5 | Rol: IS) \\
                Propiedades editables, zoom/pan, selección múltiple y persistencia 
                del grafo.
            
                \item \textbf{DT.13 – Biblioteca de recursos Cloud} 
                (25 h | Dep.: DT.12 | Rec.: R1, R4, R5 | Rol: IS) \\
                Paleta de componentes Cloud reutilizables, con cobertura inicial de AWS.
            
                \item \textbf{DT.14 – Gestión de conexiones} 
                (25 h | Dep.: DT.12 | Rec.: R1, R4, R5 | Rol: IS) \\
                Aristas del grafo y validación de conexiones entre nodos, base para 
                el código IaC.
            
                \item \textbf{DT.15 – Seguimiento y documentación del sprint} 
                (10 h | Dep.: -- | Rec.: R1, R2, R3, R4 | Rol: JP, RT) \\
                Reuniones periódicas con el director, junto con el registro de 
                actividades, resultados y trazabilidad del sprint.
            
            \end{itemize}
        

            \item \textbf{Sprint 5 – Mejoras a la implementación} 
            (13 mayo 2026 → 27 mayo 2026) 85 h
            
            \begin{itemize}[leftmargin=1.2em, itemsep=2pt]
            
                \item \textbf{DT.16 – Generador Terraform y UI} 
                (25 h | Dep.: DT.10, DT.14 | Rec.: R1, R4, R5 | Rol: IS) \\
                Motor en Rust para traducir el grafo a HCL y previsualización en UI.
            
                \item \textbf{DT.17 – Integración y empaquetado} 
                (25 h | Dep.: DT.16 | Rec.: R1, R4, R5 | Rol: IS) \\
                Persistencia de proyectos, exportación IaC y empaquetado en Linux.
            
                \item \textbf{DT.18 – Validación y seguridad} 
                (25 h | Dep.: DT.16 | Rec.: R1, R4, R5 | Rol: IS) \\
                Pruebas funcionales y análisis con Trivy.
            
                \item \textbf{DT.19 – Seguimiento y documentación del sprint} 
                (10 h | Dep.: -- | Rec.: R1, R2, R3, R4 | Rol: JP, RT) \\
                Reuniones periódicas con el director, junto con el registro de 
                actividades, resultados y trazabilidad del sprint.
            
            \end{itemize}

               \item \textbf{Sprint 6 – Memoria y preparación de la defensa} 
                (27 mayo 2026 → 15 junio 2026) 80 h
                \begin{itemize}[leftmargin=1.2em, itemsep=2pt]
                
                    \item \textbf{DT.20 – Memoria} 
                    (55 h | Dep.: DT.19 | Rec.: R1, R2, R3, R4 | Rol: RT) \\
                    Redacción y revisión final del TFG.
                
                    \item \textbf{DT.21 – Preparación de la defensa oral} 
                    (15 h | Dep.: DT.20 | Rec.: R1, R4 | Rol: JP, RT) \\
                    Elaboración de diapositivas, preparación de la demostración del 
                    prototipo y ensayo de la exposición final.

                    \item \textbf{DT.22 – Seguimiento del sprint} 
                    (10 h | Dep.: -- | Rec.: R1, R2, R3, R4 | Rol: JP, RT) \\
                    Reuniones periódicas con el director.
                
                \end{itemize}

        \end{itemize}

    \subsection{Tabla resumen de tareas}

    \justify
    La siguiente tabla recoge de forma consolidada todas las tareas del proyecto, 
    indicando para cada una su identificador, duración estimada, dependencias, 
    recursos materiales y rol humano responsable.

    \begin{table}[H]
    \centering
    \small
    \setstretch{1.0}
    \begin{tabular}{|l|l|r|l|l|l|}
    \hline
    \textbf{ID} & \textbf{Tarea} & \textbf{Horas} & \textbf{Dep.} & \textbf{Recursos} & \textbf{Rol} \\
    \hline
    GEP  & Gestión del Proyecto        & 90 & --                & R1, R2, R3        & JP, RT \\
    \hline
    DT.1  & Investigación y estado del arte     & 15 & --         & R1                & IS \\
    DT.2  & Definición de requisitos            & 25 & DT.1       & R1, R3            & JP, IS \\
    DT.3  & Selección tecnológica               & 10 & DT.1       & R1                & IS \\
    DT.4  & Seguimiento sprint 1                & 10 & --         & R1, R2, R3, R4    & JP, RT \\
    \hline
    DT.5  & Aprendizaje Rust, Tauri y React     & 30 & DT.3       & R1, R5            & IS \\
    DT.6  & Diseño de la arquitectura           & 15 & DT.5       & R1, R5            & IS \\
    DT.7  & Configuración del entorno           & 10 & DT.3       & R1, R5            & IS \\
    DT.8  & Seguimiento sprint 2                & 10 & --         & R1, R2, R3, R4    & JP, RT \\
    \hline
    DT.9  & Canvas básico                       & 45 & DT.6, DT.7 & R1, R4, R5        & IS \\
    DT.10 & Pruebas de traducción a IaC         & 15 & DT.9       & R1, R4, R5        & IS \\
    DT.11 & Seguimiento sprint 3                & 10 & --         & R1, R2, R3, R4    & JP, RT \\
    \hline
    DT.12 & Editor visual completo              & 30 & DT.9       & R1, R4, R5        & IS \\
    DT.13 & Biblioteca de recursos Cloud        & 25 & DT.12      & R1, R4, R5        & IS \\
    DT.14 & Gestión de conexiones               & 25 & DT.12      & R1, R4, R5        & IS \\
    DT.15 & Seguimiento sprint 4                & 10 & --         & R1, R2, R3, R4    & JP, RT \\
    \hline
    DT.16 & Generador Terraform y UI            & 25 & DT.10, DT.14 & R1, R4, R5      & IS \\
    DT.17 & Integración y empaquetado           & 25 & DT.16      & R1, R4, R5        & IS \\
    DT.18 & Validación y seguridad              & 25 & DT.16      & R1, R4, R5        & IS \\
    DT.19 & Seguimiento sprint 5                & 10 & --         & R1, R2, R3, R4    & JP, RT \\
    \hline
    DT.20 & Memoria                             & 55 & DT.19 & R1, R2, R3, R4    & RT \\
    DT.21 & Preparación de la defensa           & 15 & DT.20      & R1, R4            & JP, RT \\
    DT.22 & Seguimiento sprint 6                & 10 & --         & R1, R2, R3, R4    & JP, RT \\
    \hline
    \multicolumn{2}{|l|}{\textbf{Total}} & \textbf{540} & & & \\
    \hline
    \end{tabular}
    \caption{Tabla resumen de tareas del proyecto}
    \footnotesize Fuente: elaboración propia
    \label{tab:tareas}
    \end{table}

    \par



       \subsection{Gestión del riesgo}
       \label{sec:riesgos}
        \setstretch{1.2}            
        \setlength{\parskip}{1em}
        \justify
        Durante el desarrollo del proyecto pueden surgir obstáculos que afecten a su progreso. 
        Por ello, en la planificación se ha reservado un margen adicional de horas en la fase 
        de desarrollo técnico, y se han definido planes de contingencia concretos para cada 
        riesgo identificado. Los posibles riesgos se han descrito en la Sección 3.4. A 
        continuación, se indica su nivel de impacto en el proyecto y las alternativas previstas 
        para mitigarlos.

        \paragraph{R.1 – Complejidad de traducción visual → IaC}
        \textit{Probabilidad: alta | Impacto: alto}
        \justify
        Dificultad para definir una abstracción común entre el modelo visual y los dialectos 
        IaC de distintos proveedores (Terraform HCL, Pulumi).

        \textbf{Plan B:}
        \begin{itemize}[leftmargin=*, itemsep=1pt]
            \item \textbf{Tarea alternativa:} Reducir DT.13 (Biblioteca de recursos Cloud) 
            limitándola a un subconjunto mínimo de recursos AWS con plantillas HCL fijas, 
            eliminando el soporte multi-proveedor del MVP y documentándolo como trabajo futuro.
            \item \textbf{Impacto en duración:} DT.13 se reduce de 25 h a 15 h, liberando 
            10 h que se redistribuyen en DT.10 (+5 h para consolidar las pruebas de 
            traducción) y DT.16 (+5 h para el generador). La duración total del proyecto 
            no se ve afectada.
            \item \textbf{Recursos adicionales:} Ninguno. La reducción de alcance compensa 
            el esfuerzo sin requerir herramientas ni tiempo extra.
        \end{itemize}

        \paragraph{R.2 – Curva de aprendizaje de Rust}
        \textit{Probabilidad: media | Impacto: medio}
        \justify
        El aprendizaje de Rust (ownership, borrowing, lifetimes) puede retrasar el inicio 
        de la implementación técnica si la curva resulta más pronunciada de lo estimado.

        \textbf{Plan B:}
        \begin{itemize}[leftmargin=*, itemsep=1pt]
            \item \textbf{Tarea alternativa:} Ampliar DT.5 (Aprendizaje de Rust, Tauri y 
            React) en 10 h adicionales, compensadas reduciendo DT.9 (Canvas básico) en 5 h, 
            simplificando la UI inicial a un prototipo de baja fidelidad sin drag-and-drop.
            \item \textbf{Impacto en duración:} El Sprint 2 se alarga en 5 h netas, 
            absorbibles dentro del margen de contingencia reservado. La fecha de entrega 
            final no se ve comprometida.
            \item \textbf{Recursos adicionales:} Acceso a documentación oficial de Rust 
            (\textit{The Rust Programming Language}) y recursos online gratuitos (R1 
            existente). No se requieren recursos materiales adicionales.
        \end{itemize}

        \paragraph{R.3 – Estimaciones insuficientes por imprevistos}
        \textit{Probabilidad: media | Impacto: medio}
        \justify
        Imprevistos personales o académicos pueden reducir la disponibilidad en semanas 
        clave del desarrollo.

        \textbf{Plan B:}
        \begin{itemize}[leftmargin=*, itemsep=1pt]
            \item \textbf{Tarea alternativa:} Activar el margen de contingencia del Sprint 6, 
            diseñado explícitamente para absorber retrasos de hasta una semana. Si el retraso 
            supera ese margen, se reducirá el alcance de DT.21 (Preparación de la defensa) 
            de 15 h a 8 h, limitando la demo del prototipo a un caso de uso representativo.
            \item \textbf{Impacto en duración:} Un retraso de hasta 7 días se absorbe sin 
            modificar la fecha de entrega del 15 de junio. Un retraso mayor obligaría a 
            negociar una prórroga con el director, impactando únicamente el Sprint 6.
            \item \textbf{Recursos adicionales:} Ninguno. La gestión del riesgo se basa en 
            reordenación de prioridades, no en recursos extra.
        \end{itemize}
            
        
        \justify
        La distribución temporal de todas las tareas, incluyendo los márgenes de 
        contingencia descritos, se visualiza en el diagrama de Gantt de la sección 
        siguiente.

        \newpage
        \begin{landscape}  % Inicia la página en orientación horizontal
            \thispagestyle{empty} % Elimina el encabezado y pie de página SOLO en esta página
            
            \subsection{Diagrama de Gantt}
            \label{sec:gantt}
            \setstretch{1.0}            
            \setlength{\parskip}{0.5em}
            \justify
            El siguiente diagrama de Gantt refleja la planificación descrita en las 
            secciones anteriores, mostrando la secuencia de sprints, las dependencias 
            entre tareas y el margen reservado para la gestión de riesgos.
            \begin{figure}[H]
                \centering
                % \makebox permite que la imagen ignore los márgenes laterales de forma segura.
                % width=1.15\textwidth hace que la imagen sea un 15% más ancha que el texto normal. 
                % (Puedes subir este número a 1.2 o 1.25 si aún la quieres más grande).
                \makebox[\linewidth][c]{\includegraphics[width=0.75\linewidth]{TFG - Cloud IaC Visual Editor.pdf}}
                \caption{Diagrama de Gantt}
                \footnotesize Fuente: elaboración propia con ProjectLibre
            \end{figure}
            \centerline{\thepage}
        \end{landscape}  % Fin del entorno landscape
    }

    
     \newpage
    \section{Gestión económica}
   \normalsize
    \setstretch{1}
    \setlength{\parskip}{0.6em}
    \justify

    En esta sección se identifican, estiman y cuantifican los costes del proyecto.
    Siguiendo la estructura recomendada, el presupuesto se divide en: \textit{Costes 
    de Personal por Actividad (CPA)}, \textit{Costes Genéricos (CG)}, 
    \textit{Contingencias} y \textit{Costes de Imprevistos}. Todos los costes de 
    personal se derivan directamente de las tareas y horas definidas en el diagrama 
    de Gantt (Sección~\ref{sec:gantt}), garantizando la coherencia entre la 
    planificación temporal y el presupuesto.

    \subsection{Costes de Personal por Actividad (CPA)}
    \normalsize
    \setstretch{1}
    \setlength{\parskip}{0.6em}

    \subsubsection{Recursos humanos}
    \normalsize
    \setstretch{1}
    \setlength{\parskip}{0.6em}

    \justify
    Para calcular el coste de los recursos humanos, se han asignado tarifas estándar 
    de mercado a los tres roles definidos en el proyecto, consultadas a través del 
    portal de empleo \cite{noauthor_calcular_nodate}. Al salario bruto estimado por 
    hora se le ha aplicado un factor multiplicador de \textbf{1,3} para incluir los 
    costes asociados a la Seguridad Social (aproximadamente un 30\% del salario bruto, 
    correspondiente a las cotizaciones empresariales en España).

    \begin{table}[H]
    \centering
    \begin{tabular}{lcc}
    \hline
    \textbf{Rol} & \textbf{Salario bruto (\euro/h)} & \textbf{Coste con SS (\euro/h)} \\
    \hline
    Jefe de proyecto (JP) & 24,00 & 31,20 \\
    Ingeniero de software (IS) & 17,00 & 22,10 \\
    Redactor técnico (RT) & 14,00 & 18,20 \\
    \hline
    \end{tabular}
    \caption{Tarifas horarias de los recursos humanos del proyecto}
    \label{tab:costes_personal}
    \vspace{0.2cm}
    \footnotesize Fuente: elaboración propia a partir del portal de empleo citado.
    \end{table}

    \justify
    Cuando una tarea es compartida por dos roles, el coste se calcula como la media 
    aritmética de los costes horarios de ambos roles, multiplicada por el número de 
    horas de la tarea. Por ejemplo, para una tarea de 10 h compartida entre JP 
    (31,20 \euro/h) y RT (18,20 \euro/h):
    \[
    \frac{31{,}20 + 18{,}20}{2} \times 10 = 24{,}70 \times 10 = 247{,}00 \, \text{\euro}
    \]

    \justify
    Las tareas y sus costes se derivan directamente de la planificación del Gantt 
    (Sección~\ref{sec:gantt}). Para cada tarea se indica el rol responsable según 
    lo definido en la Tabla~\ref{tab:tareas}, las horas estimadas y el coste total:

    \begin{table}[H]
    \centering
    \scriptsize
    \begin{tabular}{llccc}
    \toprule
    \textbf{Sprint} & \textbf{Tarea} & \textbf{Rol} & \textbf{Horas} & \textbf{Coste (\euro)} \\
    \midrule
    
    \textbf{GEP} & Gestión del Proyecto & JP, RT & 90 & 2.223,00 \\

    \midrule

    \textbf{Sprint 1} & DT.1 Investigación y estado del arte & IS & 15 & 331,50 \\
    & DT.2 Definición de requisitos & JP, IS & 25 & 666,25 \\
    & DT.3 Selección tecnológica & IS & 10 & 221,00 \\
    & DT.4 Documentación sprint & JP, RT & 10 & 247,00 \\
    
    \midrule
    
    \textbf{Sprint 2} & DT.5 Aprendizaje tecnologías & IS & 30 & 663,00 \\
    & DT.6 Diseño arquitectura & IS & 15 & 331,50 \\
    & DT.7 Configuración entorno & IS & 10 & 221,00 \\
    & DT.8 Documentación sprint & JP, RT & 10 & 247,00 \\
    
    \midrule
    
    \textbf{Sprint 3} & DT.9 Canvas básico & IS & 45 & 994,50 \\
    & DT.10 Pruebas traducción IaC & IS & 15 & 331,50 \\
    & DT.11 Documentación sprint & JP, RT & 10 & 247,00 \\
    
    \midrule
    
    \textbf{Sprint 4} & DT.12 Editor visual & IS & 30 & 663,00 \\
    & DT.13 Biblioteca cloud & IS & 25 & 552,50 \\
    & DT.14 Gestión conexiones & IS & 25 & 552,50 \\
    & DT.15 Documentación sprint & JP, RT & 10 & 247,00 \\
    
    \midrule
    
    \textbf{Sprint 5} & DT.16 Generador Terraform & IS & 25 & 552,50 \\
    & DT.17 Integración sistema & IS & 25 & 552,50 \\
    & DT.18 Validación seguridad & IS & 25 & 552,50 \\
    & DT.19 Documentación sprint & JP, RT & 10 & 247,00 \\
    
    \midrule
    
    \textbf{Sprint 6} & DT.20 Redacción memoria & RT & 55 & 1.001,00 \\
    & DT.21 Preparación defensa & JP, RT & 15 & 370,50 \\
    & DT.22 Seguimiento sprint & JP, RT & 10 & 247,00 \\
    
    \midrule
    \multicolumn{4}{r}{\textbf{Coste total de personal (CPA)}} & \textbf{12.262,25} \euro \\
    \bottomrule
    
    \end{tabular}
    \caption{Coste estimado de recursos humanos por tarea y sprint}
    \label{tab:coste_sprints}
    \footnotesize
    El coste por hora incluye el incremento del 30\% por cotizaciones a la Seguridad Social.
    Las tareas compartidas se calculan como la media de los costes horarios de los roles implicados.
    Fuente: elaboración propia a partir del Gantt (Sección~\ref{sec:gantt}).
    \end{table}

    \subsection{Costes Genéricos (CG)}
    \normalsize
    \setstretch{1}
    \setlength{\parskip}{0.6em}
    \justify
    Los costes genéricos son aquellos que no dependen directamente del número de 
    tareas realizadas, sino de factores estructurales del proyecto: hardware, software, 
    espacio de trabajo y electricidad. Se han estimado considerando la duración total 
    del proyecto de \textbf{540 horas} (450 h de desarrollo técnico y 90 h de gestión), 
    distribuidas a lo largo de \textbf{5 meses} (febrero–junio 2026).

    \subsubsection{Hardware}
    \normalsize
    \setstretch{1}
    \setlength{\parskip}{0.6em}
    \justify
    En un entorno profesional, el hardware sería compartido entre varios miembros del 
    equipo. En este proyecto, sin embargo, todos los roles son desempeñados por una 
    única persona, por lo que se contabiliza un único puesto de trabajo. El coste 
    imputable al proyecto no es el precio de adquisición completo, sino la parte 
    proporcional correspondiente a las horas de uso durante el proyecto, calculada 
    mediante amortización lineal.

    \begin{table}[H]
    \centering
    \begin{tabular}{lc}
    \hline
    \textbf{Dispositivo} & \textbf{Coste de adquisición (\euro)} \\
    \hline
    Torre PC (AMD Ryzen 7 5800X, RTX 3060, 32 GB RAM, 1 TB NVMe) & 1.300,00 \\
    Monitor HP Omen & 125,10 \\
    Monitor Dell & 110,00 \\
    Teclado & 23,99 \\
    Ratón & 5,99 \\
    \hline
    \textbf{Total adquisición} & \textbf{1.565,08} \euro \\
    \hline
    \end{tabular}
    \caption{Coste de adquisición del hardware utilizado}
    \label{tab:hardware_coste}
    \footnotesize Fuente: elaboración propia.
    \end{table}

    \justify
    \textbf{Amortización:} Se considera una vida útil media de \textbf{6 años}, 
    criterio habitual para equipos informáticos en entornos profesionales. Con 
    251 días laborables al año y 8 horas por jornada, el total de horas de vida 
    útil es:
    \[
    251 \times 8 \times 6 = 12.048 \, \text{h}
    \]
    El coste amortizado por hora de uso es:
    \[
    \frac{1.565{,}08}{12.048} = 0{,}1299 \, \text{\euro/h}
    \]
    Para las 540 horas del proyecto:
    \[
    0{,}1299 \times 540 = \mathbf{70{,}15} \, \text{\euro}
    \]

    \subsubsection{Software}
    \normalsize
    \setstretch{1}
    \setlength{\parskip}{0.6em}
    \justify
    Todas las herramientas empleadas (Rust, Tauri, React, GitHub, Trello, \LaTeX{}, 
    VS Code) son gratuitas y de código abierto, por lo que el coste de licencias es 
    \textbf{0,00 \euro}. Esta decisión es coherente con el enfoque open source del 
    proyecto y elimina cualquier dependencia de software propietario.

    \subsubsection{Espacio de trabajo}
    \normalsize
    \setstretch{1}
    \setlength{\parskip}{0.6em}
    \justify
    El proyecto se desarrolla en el domicilio del autor. Para estimar el coste 
    imputable al proyecto se consideran dos componentes: el alquiler mensual 
    proporcional y otros gastos generales fijos durante la duración del proyecto.

    \begin{itemize}[leftmargin=*, itemsep=1pt]
    \item \textbf{Alquiler y suministros}: 2.030 \euro/mes × 5 meses = 10.150 \euro
    \item \textbf{Otros gastos generales} (limpieza, material de oficina, etc.): 
    2.000 \euro\ fijos para toda la duración
    \end{itemize}

    \[
    2.030 \times 5 + 2.000 = \mathbf{12.150} \, \text{\euro}
    \]

    \subsubsection{Electricidad}
    \normalsize
    \setstretch{1}
    \setlength{\parskip}{0.6em}
    \justify
    El coste de electricidad se calcula en dos partes: el consumo de los dispositivos 
    del puesto de trabajo durante las 540 horas del proyecto, y el consumo general 
    del domicilio durante los 5 meses.

    \justify
    Para el consumo de dispositivos, se estima la potencia media de cada equipo y 
    se multiplica por las 540 horas totales del proyecto:

    \begin{table}[H]
    \centering
    \begin{tabular}{lccc}
    \hline
    \textbf{Dispositivo} & \textbf{Potencia (W)} & \textbf{Horas uso} & \textbf{Consumo (kWh)} \\
    \hline
    Torre PC & 250 & 540 & 135,00 \\
    Monitor HP Omen & 35 & 540 & 18,90 \\
    Monitor Dell & 30 & 540 & 16,20 \\
    Teclado & 5 & 540 & 2,70 \\
    Ratón & 3 & 540 & 1,62 \\
    \hline
    \textbf{Total consumo dispositivos} & & & \textbf{174,42 kWh} \\
    \hline
    \end{tabular}
    \caption{Consumo eléctrico estimado de los dispositivos del proyecto}
    \label{tab:consumo_electrico}
    \footnotesize Fuente: elaboración propia.
    \end{table}

    \justify
    Según los datos oficiales del precio medio de la electricidad en España 
    \cite{noauthor_precio_nodate}, el coste medio es de \textbf{0,1975 \euro/kWh}. 
    El coste de los dispositivos es:
    \[
    174{,}42 \times 0{,}1975 \approx 34{,}45 \, \text{\euro}
    \]

    \justify
    A esto se añade el consumo proporcional del domicilio durante los 5 meses del 
    proyecto, estimado en 60 \euro/mes basándose en la factura media de un hogar 
    español con equipos ofimáticos:
    \[
    60 \times 5 = 300 \, \text{\euro}
    \]
    \[
    34{,}45 + 300 \approx \mathbf{334{,}45} \, \text{\euro}
    \]

    \subsubsection{Resumen de costes genéricos}
    \normalsize
    \setstretch{1}
    \setlength{\parskip}{0.6em}

    \begin{table}[H]
    \centering
    \begin{tabular}{lc}
    \hline
    \textbf{Concepto} & \textbf{Coste (\euro)} \\
    \hline
    Hardware (amortización) & 70,15 \\
    Software & 0,00 \\
    Espacio de trabajo & 12.150,00 \\
    Electricidad & 334,45 \\
    \hline
    \textbf{Total costes genéricos (CG)} & \textbf{12.554,60} \\
    \hline
    \end{tabular}
    \caption{Resumen de costes genéricos del proyecto}
    \label{tab:costes_genericos}
    \footnotesize Fuente: elaboración propia.
    \end{table}

    \subsection{Contingencia}
    \normalsize
    \setstretch{1}
    \setlength{\parskip}{0.6em}
    \justify
    En proyectos informáticos es habitual reservar un margen de contingencia para 
    cubrir desviaciones en los costes estimados. Para este proyecto se aplica un 
    \textbf{15\%} sobre la suma de CPA y CG, porcentaje estándar en proyectos de 
    desarrollo software de alcance medio. Este 
    margen cubre posibles incrementos en horas de trabajo debidos a imprevistos 
    técnicos no contemplados en los riesgos identificados.

    \[
    \text{Contingencia} = (CPA + CG) \times 0{,}15 = 
    (12.262{,}25 + 12.554{,}60) \times 0{,}15 \approx \mathbf{3.722{,}53} \, \text{\euro}
    \]

    \subsection{Costes de Imprevistos}
    \normalsize
    \setstretch{1}
    \setlength{\parskip}{0.6em}
    \justify
    Los imprevistos se derivan directamente de los riesgos identificados en la 
    Sección~\ref{sec:riesgos}. Siguiendo la metodología GEP, cada imprevisto no 
    se costea al 100\% sino según la probabilidad estimada de que ocurra el riesgo: 
    para un riesgo con probabilidad $p$, se imputa $p \times \text{coste\_plan\_B}$. 
    Esto configura un sistema de compensación: si solo ocurre uno de los imprevistos, 
    lo presupuestado por los que no ocurren cubre parcialmente el exceso del que sí ocurre.

    \begin{table}[H]
    \centering
    \begin{tabular}{lp{3.5cm}ccrr}
    \hline
    \textbf{Riesgo} & \textbf{Plan B} & \textbf{Prob.} & \textbf{Horas} & 
    \textbf{Coste plan B (\euro)} & \textbf{Imprevisto (\euro)} \\
    \hline
    R.1 & Reducción DT.13, redistribución 10 h IS & 70\% & 10 h IS & 221,00 & 154,70 \\
    R.2 & Ampliación DT.5 en 5 h netas IS & 40\% & 5 h IS  & 110,50 &  44,20 \\
    R.3 & Reducción DT.21 en 7 h JP/RT & 40\% & 7 h JP/RT & 172,90 &  69,16 \\
    \hline
    \textbf{Total} & & & & & \textbf{268,06} \\
    \hline
    \end{tabular}
    \caption{Estimación de costes de imprevistos del proyecto}
    \label{tab:imprevistos}
    \footnotesize Fuente: elaboración propia a partir de la gestión de riesgos 
    (Sección~\ref{sec:riesgos}).
    \end{table}

    \subsection{Presupuesto Total}
    \normalsize
    \setstretch{1}
    \setlength{\parskip}{0.6em}
    \justify
    La siguiente tabla resume el presupuesto total del proyecto, integrando todas 
    las partidas calculadas:

    \begin{table}[H]
    \centering
    \begin{tabular}{lc}
    \hline
    \textbf{Concepto} & \textbf{Importe (\euro)} \\
    \hline
    Costes de Personal por Actividad (CPA) & 12.262,25 \\
    Costes Genéricos (CG) & 12.554,60 \\
    \hline
    \textbf{Subtotal (CPA + CG)} & \textbf{24.816,85} \\
    Contingencias (15\% sobre subtotal) & 3.722,53 \\
    Imprevistos & 268,06 \\
    \hline
    \textbf{Presupuesto total estimado} & \textbf{28.807,44} \\
    \hline
    \end{tabular}
    \caption{Resumen del presupuesto total estimado del proyecto}
    \label{tab:presupuesto_total}
    \footnotesize Fuente: elaboración propia.
    \end{table}

    \justify
    El presupuesto total estimado asciende a \textbf{28.807,44 \euro}. Las partidas 
    dominantes son los costes de personal (42,6\% del total) y el espacio de trabajo 
    (42,2\%), lo que es coherente con un proyecto de desarrollo software individual 
    realizado desde el domicilio. Los imprevistos representan únicamente el 0,9\% 
    del total, reflejo de que los planes B definidos priorizan la reorganización de 
    tareas existentes sobre la incorporación de horas adicionales. En el contexto 
    académico de un TFG, donde el estudiante no percibe remuneración, el coste 
    efectivo real se limita a los costes genéricos (\textbf{12.554,60 \euro}), 
    aunque el presupuesto refleja el valor económico real si el proyecto se 
    desarrollara en un entorno profesional.

    \subsection{Control de Gestión}
    \normalsize
    \setstretch{1}
    \setlength{\parskip}{0.6em}
    \justify
    El objetivo del control de gestión es comparar y evaluar las desviaciones entre 
    el presupuesto planificado y los costes reales en que va incurriendo el proyecto 
    a medida que avanza. Siguiendo la metodología del GEP, el control se articula 
    en torno a tres preguntas fundamentales:

    \begin{itemize}[leftmargin=*, itemsep=1pt]
    \item \textbf{¿Dónde} se ha producido la desviación? (en qué tarea o sprint)
    \item \textbf{¿Por qué} se ha producido? (desviación en tarifa o en eficiencia)
    \item \textbf{¿De qué cuantía} es la desviación?
    \end{itemize}

    \justify
    Para responder a estas preguntas se definen los siguientes indicadores numéricos, 
    que se revisarán al final de cada sprint coincidiendo con la reunión de seguimiento 
    con el director:

    \subsubsection*{Desvío en tarifa (precio)}
    \justify
    Mide si el coste real por hora difiere del coste estimado, manteniendo constante 
    el consumo de horas:

    \[
    D_{tarifa} = (C_{h,est} - C_{h,real}) \times H_{real}
    \]

    donde $C_{h,est}$ es el coste horario estimado del rol, $C_{h,real}$ el coste 
    horario real y $H_{real}$ las horas realmente consumidas. En este proyecto, al 
    ser el mismo autor quien desempeña todos los roles, este desvío sería nulo salvo 
    que se requiriera subcontratar trabajo externo.

    \subsubsection*{Desvío en eficiencia (consumo)}
    \justify
    Mide si se han consumido más o menos horas de las planificadas para una tarea, 
    aplicando el coste estimado como referencia:

    \[
    D_{eficiencia} = (H_{est} - H_{real}) \times C_{h,est}
    \]

    donde $H_{est}$ son las horas planificadas para la tarea. Un valor negativo indica 
    sobreconsumo de horas. Este es el indicador principal del proyecto, ya que el 
    principal riesgo de desviación es la infraestimación de horas de desarrollo.

    \subsubsection*{Desvío total de mano de obra por sprint}
    \justify
    Integra ambos efectos para cada sprint $k$:

    \[
    D_{total,k} = \sum_{i \in Sprint_k} \left[ (C_{h,est,i} - C_{h,real,i}) \times 
    H_{real,i} + (H_{est,i} - H_{real,i}) \times C_{h,est,i} \right]
    \]

    \justify
    La siguiente tabla muestra los valores de referencia planificados por sprint, 
    que servirán como línea base para el cálculo de los desvíos reales:

    \begin{table}[H]
    \centering
    \small
    \begin{tabular}{lrrrr}
    \hline
    \textbf{Sprint} & \textbf{Horas planif.} & \textbf{Coste planif. (\euro)} & 
    \textbf{Coste acum. (\euro)} & \textbf{Umbral alerta (5\%)} \\
    \hline
    GEP         & 90  & 2.223,00 & 2.223,00  & 111,15 \\
    Sprint 1    & 60  & 1.465,75 & 3.688,75  & 73,29  \\
    Sprint 2    & 65  & 1.462,50 & 5.151,25  & 73,13  \\
    Sprint 3    & 70  & 1.573,00 & 6.724,25  & 78,65  \\
    Sprint 4    & 90  & 2.015,00 & 8.739,25  & 100,75 \\
    Sprint 5    & 85  & 1.904,50 & 10.643,75 & 95,23  \\
    Sprint 6    & 80  & 1.618,50 & 12.262,25 & 80,93  \\
    \hline
    \textbf{Total} & \textbf{540} & \textbf{12.262,25} & & \\
    \hline
    \end{tabular}
    \caption{Línea base del presupuesto por sprint para el control de gestión}
    \label{tab:control_gestion}
    \footnotesize Fuente: elaboración propia a partir de la Tabla~\ref{tab:coste_sprints}.
    \end{table}

    \justify
    Si al final de un sprint el $D_{total,k}$ supera el \textbf{5\%} del coste 
    planificado del sprint (columna ``Umbral alerta''), se activará el siguiente 
    protocolo de actuación en tres niveles:

    \begin{enumerate}[leftmargin=*, itemsep=1pt]
    \item \textbf{Nivel 1 — Revisión de alcance:} reducir el alcance de las tareas 
    pendientes del sprint siguiente para compensar el sobrecoste.
    \item \textbf{Nivel 2 — Imprevistos:} si la desviación acumulada supera la 
    partida de imprevistos (268,06 \euro), se activará el fondo de contingencia 
    (3.722,53 \euro).
    \item \textbf{Nivel 3 — Renegociación:} si el coste total proyectado supera el 
    presupuesto total más contingencias, se negociará con el director una reducción 
    formal del alcance del MVP.
    \end{enumerate}

    \justify
    Este sistema implementa el concepto de \textbf{presupuesto flexible} descrito 
    en la metodología GEP: el presupuesto pactado es uno, pero a partir del control 
    por sprint pueden generarse estimaciones actualizadas del coste final del proyecto, 
    permitiendo anticipar desviaciones antes de que sean irrecuperables.


    \newpage
    \section{Informe de Sostenibilidad}
    \normalsize
    \setstretch{1}
    \setlength{\parskip}{0.6em}
    \justify
     Se evalúan los ámbitos ambiental, económico y social del proyecto, incluyendo una autoevaluación de competencias en sostenibilidad.
    
    \subsection{Autoevaluación}
    \normalsize
    \setstretch{1}
    \setlength{\parskip}{0.6em}
    \justify
    Tras completar el cuestionario EDINSOST de autoevaluación de competencias en
    sostenibilidad, he identificado puntos fuertes y áreas de mejora en los tres
    ámbitos de la sostenibilidad: ambiental, social y económica. En ingeniería se
    tiende a priorizar la dimensión tecnológica y económica, relegando el impacto
    ambiental y social. El cuestionario me ha permitido incorporar conceptos que
    no tenía plenamente asimilados, especialmente en economía circular, economía
    social y economía ecológica.
    
    \justify
    En cuanto a mis puntos fuertes, tengo un conocimiento razonable de la
    sostenibilidad ambiental en el sector tecnológico, especialmente en consumo
    energético de infraestructuras Cloud y eficiencia computacional. Soy consciente
    de que las TIC contribuyen significativamente a las emisiones globales de CO2 y
    de que soluciones más eficientes generan un impacto ambiental positivo. Además,
    mi formación en gestión económica me permite estimar costes y valorar la
    viabilidad de un proyecto con cierta solvencia.
    
    \justify
    Respecto a las competencias adquiridas, he aprendido a considerar la huella
    ecológica de productos a lo largo de todo su ciclo de vida, aplicando
    estrategias de reducción, reutilización y reciclaje. En el ámbito social, he
    reflexionado sobre el impacto de mis proyectos en la salud, seguridad y
    justicia social, incluyendo accesibilidad, diversidad y equidad.
    
    \justify
    En cuanto a las áreas de mejora, mi conocimiento sobre métricas cuantitativas
    de sostenibilidad, como el análisis del ciclo de vida o las metodologías de
    huella de carbono, es aún limitado. Tampoco he trabajado con marcos formales
    de sostenibilidad social, aunque comprendo sus principios fundamentales.
    
    \justify
    En conclusión, la autoevaluación ha reforzado la necesidad de integrar la
    sostenibilidad de forma transversal desde las etapas más tempranas del
    proyecto, y reconozco que debo seguir profundizando en estas competencias
    para un análisis más avanzado.
    

    \subsubsection{Dimensión Económica}
    \normalsize
    \setstretch{1}
    \setlength{\parskip}{0.6em}
    \justify

    En relación con el Proyecto Puesto en Producción (PPP), el coste total estimado
    asciende a \textbf{28.807,44} \euro, como se muestra en la Tabla
    \ref{tab:presupuesto_total}. Este importe incluye los Costes de Personal por
    Actividad (CPA), los Costes Genéricos (CG), un margen de contingencia del 15\%
    y los costes de imprevistos, calculados según la probabilidad estimada de 
    materialización de cada riesgo identificado. En el contexto académico individual, 
    el coste efectivo se limita a los recursos materiales y amortizaciones
    (\textbf{12.554,60} \euro), aunque el presupuesto refleja el valor económico
    real de los recursos implicados.

    \justify
    Respecto al estado del arte, herramientas comerciales de modelado visual de IaC
    como \textit{Brainboard} tienen costes entre
    100 \euro{} y 500 \euro{} mensuales por usuario. La solución propuesta, al ser
    de código abierto, elimina estos costes y supone una alternativa económica
    relevante para equipos pequeños, startups y entornos educativos. Además, su
    interfaz visual reduce la curva de aprendizaje de IaC, disminuyendo el tiempo
    de incorporación de nuevos perfiles técnicos y generando un ahorro indirecto
    en formación y productividad.
    
    
    \subsubsection{Dimensión Ambiental}
    \normalsize
    \setstretch{1}
    \setlength{\parskip}{0.6em}
    \justify

    \justify
    \textbf{Impacto de la realización (PPP):} El impacto ambiental directo de la
    ejecución de este TFG es reducido y se limita principalmente al consumo
    eléctrico de los dispositivos durante las 540 horas de trabajo, estimado en
    \textbf{34,45} \euro{}. Además, no se ha adquirido nuevo hardware, reutilizando
    equipos existentes y minimizando el impacto asociado a su fabricación y
    transporte.

    \justify
    \textbf{Minimización del impacto (PPP):} Se ha evitado el uso de servidores
    físicos dedicados, trabajando en entornos locales con herramientas de software
    libre. El desarrollo individual en un único puesto de trabajo implica un coste
    de \textbf{70,15} \euro{} por amortización de hardware y
    \textbf{334,45} \euro{} en consumo eléctrico total, sin necesidad de
    infraestructura Cloud durante esta fase.

    \justify
    \textbf{Vida útil y mejora de la solución:} La solución presenta un potencial
    impacto ambiental positivo en fase operativa. Su interfaz visual facilita la
    gestión eficiente de infraestructuras Cloud, ayudando a evitar
    sobreaprovisionamientos y a identificar recursos ociosos, lo que favorece
    prácticas de eficiencia energética en centros de datos.
    
    
    \subsubsection{Dimensión Social}
    \normalsize
    \setstretch{1.1}
    \setlength{\parskip}{0.6em}
    \justify
    
    \textbf{Aporte a nivel personal (PPP):} La realización de este proyecto supone
    un crecimiento profesional significativo al asumir todos los roles del
    proyecto (jefe de proyecto, ingeniero de software y redactor técnico) durante
    540 horas. Esto ha implicado aplicar metodologías ágiles como Scrum y Kanban,
    trabajar con tecnologías como Rust y Node.js, y consolidar competencias
    técnicas junto con habilidades de autogestión.

    \justify
    \textbf{Estado del arte y mejora social:} El acceso a herramientas visuales de
    IaC suele estar limitado a soluciones comerciales o a perfiles con experiencia
    avanzada en Terraform o Pulumi, generando barreras para estudiantes, perfiles
    junior y pequeñas organizaciones. La solución propuesta, al ser de código
    abierto y visual, reduce estas barreras y favorece la incorporación de nuevos
    perfiles al ámbito del Cloud Computing.

    \justify
    \textbf{Necesidad real del proyecto:} Los equipos Cloud afrontan problemas de
    interoperabilidad y complejidad entre plataformas como AWS, GCP y Azure,
    especialmente para perfiles con menor experiencia. La herramienta propuesta
    ofrece una alternativa accesible e interoperable alineada con la filosofía
    Open Source, contribuyendo a reducir la brecha tecnológica.


    \newpage
    \section{Arquitectura del sistema}
    {
        \setstretch{1.2}
        \setlength{\parskip}{1em}

        \justify
        Esta sección describe la arquitectura general de la solución desde una perspectiva de visión global. El
        objetivo es mostrar cómo las decisiones tomadas para poder cumplir con los objetivos y requisitos definidos en la (Sección~\ref{sec:Alcance})

        \subsection{De los requisitos a la arquitectura}
        \justify
        Los subobjetivos definidos en la Sección~\ref{sec:Alcance_subobjetivos} imponen varias
        restricciones sobre la forma del sistema. Se requiere una interfaz gráfica
        capaz de representar componentes Cloud reutilizables y conectables; un
        mecanismo que traduzca de forma fiable ese diseño visual a código de
        \textbf{Infraestructura como Código (IaC)}; una arquitectura modular y
        extensible que permita incorporar nuevos proveedores en el futuro; y, como
        requisitos transversales, un enfoque \textbf{open source} y un coste de
        adopción reducido para el usuario final.

        \justify
        Estas restricciones descartan algunas opciones como: una aplicación puramente web, alojada en un servidor remoto,
        introduciría una dependencia de infraestructura externa y complicaría tanto
        la instalación local como la edición de ficheros del usuario en disco. Una
        aplicación de escritorio en lenguaje propietario, comprometería
        la portabilidad y la reutilización del código. La opción que ha resultado
        coherente con el conjunto de requisitos ha sido la de una \textbf{aplicación
        de escritorio multiplataforma}, que opera íntegramente en local y se apoya en herramientas estándar del ecosistema
        Cloud para su funcionamiento.

        \subsection{Visión global del sistema}
        \justify
        El sistema se organiza en torno a cuatro bloques principales que colaboran
        mediante interfaces estables: una \textbf{capa de presentación}, una
        \textbf{capa de aplicación}, una \textbf{capa de infraestructura} y una
        \textbf{pipeline de extracción de recursos}. Esta organización permite separar la
        interacción con el usuario, la lógica del dominio, la integración con el
        entorno local y la obtención del catálogo de recursos.
        El núcleo semántico que vertebra estos bloques es el 
        \textbf{modelo de dominio}, cuya estructura se detalla en la 
        Sección~\ref{sec:modelo_dominio}.


        \justify
        La \textbf{pipeline de extracción de recursos} obtiene el catálogo de recursos
        disponibles para cada proveedor cloud a partir de la información expuesta por
        Terraform. Esta pieza garantiza la \textbf{coherencia} y la
        \textbf{actualización} del catálogo, ya que se apoya en la fuente de verdad
        oficial del proveedor y evita mantener descripciones manuales susceptibles de
        desfase.

        \begin{figure}[H]
            \centering
            \includegraphics[
                width=0.75\textwidth,
                height=10cm,
                keepaspectratio,
                trim=0 0.5cm 0 0,
                clip
            ]{Pipeline.pdf}
            \caption{Pipeline de extracción de recursos}
            \footnotesize Fuente: elaboración propia
            \label{fig:pipeline}
        \end{figure} 
        

        \justify
        La \textbf{capa de presentación} es responsable de la interacción con el
        usuario: presenta el lienzo, el catálogo de recursos, los formularios de
        edición y los mecanismos de navegación. Está construida sobre un stack web
        moderno, lo que facilita el desarrollo mediante componentes reutilizables y
        permite una posible evolución futura hacia una versión web sin reescritura
        sustancial de esta capa.

        \justify
        La \textbf{capa de aplicación} concentra la lógica funcional del sistema:
        gestiona el estado del proyecto, valida operaciones, coordina los distintos
        componentes y traduce el diseño interno a código \textbf{Infraestructura como
        Código (IaC)}. En esta capa reside la lógica de negocio que determina cómo se
        crean, modifican y serializan los elementos del sistema.

        \justify
        Por último, la \textbf{capa de infraestructura} conecta la aplicación con el
        entorno local. Se encarga de invocar Terraform como subproceso, gestionar las
        credenciales del proveedor cloud y acceder al sistema de ficheros del proyecto.
        Esta capa reduce el acoplamiento con herramientas externas y mejora la
        \textbf{portabilidad} de la solución.

        \justify
        Para situar el sistema en su entorno y mostrar la relación entre sus
        bloques internos se emplea la notación \textbf{C4}~\cite{noauthor_home_nodate} en dos
        niveles complementarios. El \textbf{nivel 1 (contexto)} sitúa la aplicación
        frente al usuario y los sistemas externos con los que interactúa, mientras
        que el \textbf{nivel 2 (contenedores)} descompone la aplicación en sus
        bloques internos y detalla las tecnologías asociadas a cada uno.

        \begin{figure}[H]
            \centering
            \includegraphics[
                width=1\textwidth,
                height=15cm,
                keepaspectratio
            ]{nivel1.pdf}
            \caption{Vista de contexto del sistema (C4, nivel 1)}
            \footnotesize Fuente: elaboración propia
            \label{fig:arquitectura_c4}
        \end{figure}



        \subsection{Stack tecnológico}
        \justify
        El stack tecnológico se ha seleccionado de forma que cada decisión responda a un
        requisito concreto de los descritos en la Sección~\ref{sec:Alcance}: ejecución
        \textbf{multiplataforma} y \textbf{local}, arquitectura \textbf{modular y
        extensible}, y compromiso con un enfoque \textbf{open source}. A continuación se
        justifica la elección de cada pieza.

        \subsubsection*{Tauri como contenedor de aplicación de escritorio}
        \justify
        Para empaquetar la aplicación como ejecutable de escritorio se ha optado por
        \textbf{Tauri} frente a alternativas como Electron~\cite{noauthor_build_nodate}. Tauri utiliza el motor de
        renderizado nativo del sistema operativo en lugar de embeber un Chromium completo,
        lo que reduce drásticamente el tamaño del binario y el consumo de memoria
        respecto a un equivalente en Electron. Además, su backend nativo en Rust permite
        invocar subprocesos del sistema operativo de forma controlada, condición
        necesaria para integrar la CLI de Terraform y la terminal embebida descritas en
        la capa de infraestructura.

        \subsubsection*{Rust en el backend nativo}
        \justify
        El backend nativo, responsable de la capa de infraestructura, se ha implementado
        en \textbf{Rust}. La elección se sostiene en tres motivos: la seguridad de
        memoria que ofrece su sistema de \textit{ownership} y \textit{borrowing}, lo que
        elimina toda una clase de errores en código que gestiona procesos, ficheros y
        credenciales; el rendimiento sostenido al lanzar y parsear en \textit{streaming}
        la salida de Terraform, que puede ser voluminosa; y la robustez de su ecosistema
        para concurrencia (lectores de \texttt{stdout}/\texttt{stderr} en paralelo) y
        manejo de procesos interactivos, requerida para la confirmación de \texttt{apply}
        y \texttt{destroy}.

        \subsubsection*{React con TypeScript en la capa de presentación}
        \justify
        La interfaz se construye con \textbf{React} y \textbf{TypeScript}. React aporta
        un modelo de componentes reutilizables idóneo para una UI con paneles
        independientes (catálogo, lienzo, inspector, terminal) y dispone de un
        ecosistema maduro de bibliotecas para canvas, \textit{drag-and-drop} y edición
        visual. TypeScript añade tipado estático al contrato entre el frontend y los
        comandos Tauri del backend, lo que detecta inconsistencias en tiempo de
        compilación y mejora la mantenibilidad del código a medio plazo.

        \subsubsection*{Go en la pipeline de extracción de recursos}
        \justify
        La pipeline descrita en la Sección~\ref{sec:pipeline} se ha implementado en
        \textbf{Go}~\cite{noauthor_go_nodate}. La razón principal es práctica: HashiCorp mantiene una biblioteca
        oficial en Go para interpretar el schema de los proveedores de Terraform, lo
        que evita reimplementar el parser. Además, Go es especialmente adecuado para
        herramientas de procesado masivo de ficheros JSON \cite{noauthor_json_nodate}, que es exactamente la
        naturaleza del problema. Es importante recalcar que la pipeline es un
        \textbf{preproceso offline}: produce ficheros estáticos que pasan a formar parte
        de la distribución de la aplicación, por lo que el usuario final no necesita
        tener Go instalado ni ejecutar la pipeline.

        \justify
        En conjunto, todas las piezas del stack son \textbf{open source y gratuitas}, lo
        que es coherente con el subobjetivo de licenciamiento abierto y elimina cualquier
        dependencia con software propietario.

        \begin{figure}[H]
            \centering
            \includegraphics[
                width=1.1\textwidth,
                height=15cm
            ]{nivel2.pdf}
            \caption{Vista de contenedores del sistema (C4, nivel 2)}
            \footnotesize Fuente: elaboración propia
            \label{fig:arquitectura_c4_contenedores}
        \end{figure}

        \subsection{Modelo de dominio}
        \justify
        Una decisión central de la arquitectura es situar el \textbf{modelo de dominio}
        como fuente de verdad única del sistema. Mientras que el \textbf{modelo de
        presentación} describe el estado visual, posiciones, selecciones, atributos
        gráficos. El modelo de dominio captura la semántica real de la
        infraestructura: qué recursos existen, cómo se relacionan y qué valores tienen
        sus propiedades, con independencia de cómo se representen visualmente o del
        formato IaC al que se traduzcan. Su estructura interna, ciclo de vida y esquema
        de persistencia se desarrollan en detalle en la
        Sección~\ref{sec:modelo_dominio}.


        \subsection{Síntesis}
        \justify
        En conjunto, la arquitectura propuesta separa de forma explícita la
        representación visual, la lógica de aplicación, la integración con
        infraestructura y la obtención del catálogo de recursos. El sistema se apoya en
        un \textbf{modelo de dominio} único como fuente de verdad, mientras que la capa
        de presentación actúa como proyección visual de ese estado, la capa de
        aplicación concentra las operaciones necesarias para modificarlo, validarlo y
        traducirlo a \textbf{Infraestructura como Código (IaC)} y la capa de infraestructura se encarga de la interacción con el entorno local y las herramientas externas.

        \justify
        Esta separación reduce el acoplamiento entre interfaz y lógica, mejora la
        mantenibilidad y facilita la evolución del sistema. Además, el uso de un
        \textbf{catálogo extensible} permite mantener la neutralidad respecto al
        proveedor cloud y abre la puerta a incorporar nuevos entornos sin modificar el
        núcleo funcional de la aplicación.
        \par

        
    }

    \newpage
    \section{Casos de uso}
    {
        \setstretch{1.2}
        \setlength{\parskip}{1em}
        \label{sec:casos_uso}

        \justify
        Una vez descrita la arquitectura, esta sección resume las funcionalidades
        que la aplicación pone a disposición del usuario. El actor principal es el \textbf{Usuario}, ya caracterizado en
        la Sección~\ref{sec:actores_implicados}. Los \textbf{sistemas externos}
        que intervienen en estos flujos (\textit{Terraform CLI} y
        \textit{Proveedor cloud}) son los mismos que aparecen en el diagrama de
        contexto de la Figura~\ref{fig:arquitectura_c4}.

        \subsection{Visión general de las funcionalidades}
        \justify
        La aplicación cubre el ciclo de vida completo de una arquitectura cloud,
        desde su diseño inicial hasta la comparación entre lo diseñado y lo
        realmente desplegado. Las funcionalidades pueden agruparse en cuatro
        bloques complementarios.

        \justify
        El primer bloque corresponde a la \textbf{gestión de proyectos y vistas}:
        el usuario puede crear, abrir y guardar proyectos, así como organizar
        cada proyecto en una o varias vistas (típicamente \textit{dev}, \textit{staging} o
        \textit{prod}), cada una con su propio estado y sus propios archivos. El segundo bloque agrupa las acciones de \textbf{diseño
        visual}: arrastrar recursos desde el catálogo al lienzo, editar sus
        atributos en el inspector de propiedades y conectarlos entre sí mediante
        \textit{enlaces} que generan automáticamente las referencias HCL
        correspondientes. Como
        complemento, el usuario también puede \textbf{editar el HCL
        manualmente} desde el editor embebido: el sistema detecta los cambios y
        los reconcilia con el modelo de dominio, manteniendo coherentes el
        lienzo, el inspector y el código.

        \justify
        El tercer bloque comprende las acciones de \textbf{despliegue}, que
        delegan en la \textit{Terraform CLI} a través de la capa de
        infraestructura. El usuario puede ejecutar \texttt{plan} para previsualizar
        los cambios, \texttt{apply} para materializarlos en el proveedor cloud, y
        \texttt{destroy} para revertir el despliegue. La salida del subproceso se
        transmite en tiempo real y, en el caso de
        \texttt{apply}, la confirmación interactiva se intercepta y se envía desde
        la propia interfaz, sin requerir que el usuario abandone la aplicación.
        El cuarto bloque cubre las funcionalidades de \textbf{observación}:
        la vista \textit{Diff} proyecta sobre el lienzo los cambios previstos por
        el último \texttt{plan}, mientras que la vista \textit{Cloud} consulta el
        estado real de la
        infraestructura desplegada, permitiendo identificar de un vistazo
        cualquier divergencia entre diseño y realidad.
    }

    \newpage
    \section{Pipeline de extracción de recursos}
    {
        \setstretch{1.2}
        \setlength{\parskip}{1em}
        \label{sec:pipeline}
        \justify
        Esta sección describe cómo se obtienen los recursos de un proveedor cloud. 
        Aunque el desarrollo se ha realizado sobre el
        proveedor de \textbf{Amazon Web Services} para acotar el alcance del MVP, el
        proceso aquí descrito es genérico: el mismo pipeline aplicado sobre Google
        Cloud o Microsoft Azure produciría un catálogo equivalente sin requerir
        cambios en la aplicación que lo consume.

        \subsection{Motivación: ¿por qué no codificar el catálogo a mano?}
        \justify
        Una primera aproximación al problema, sería escribir
        manualmente la lista de recursos cloud que la aplicación soporta, junto con
        sus atributos y las plantillas de código IaC asociadas. Este enfoque tiene tres
        problemas fundamentales que justifican la inversión en un pipeline automatizado.

        \justify
        En primer lugar, el catálogo de un proveedor cloud mayor es extraordinariamente
        amplio: el provider de AWS para Terraform expone aproximadamente \textbf{1.500
        recursos} y \textbf{600 fuentes de datos}. Codificarlos manualmente no es viable en el marco
        temporal de un Trabajo Final de Grado, ni siquiera para un subconjunto
        razonable.

        \justify
        En segundo lugar, esa información cambia con cada versión del
        proveedor: nuevas funcionalidades introducen nuevos atributos, otras quedan
        marcadas como obsoletas, y los tipos de algunos campos se refinan. Mantener un
        catálogo manual sincronizado con esa evolución sería una tarea de
        \textbf{mantenimiento} continuado y propensa a errores.
        
        \justify
        En tercer lugar, la documentación pública del proveedor está pensada para ser
        leída, no parseada, y omite información estructural relevante de los atributos.

        \justify
        La solución natural es apoyarse en la propia herramienta IaC.
        Terraform expone, mediante un comando de su CLI, una representación completa del schema del proveedor en formato JSON.
        Esa representación procede directamente del binario oficial del proveedor lo que garantiza que sea
        \textbf{fidedigna} (está siempre alineada con la versión del proveedor que se
        está utilizando). Como se observa en la Figura~\ref{fig:pipeline}, el proceso completo de la pipeline transforma el schema del proveedor en los artefactos utilizados por la aplicación.

        \subsection{Visión general de la pipeline}
        \justify
        La pipeline está organizada en cuatro pasos encadenados. El primero genera el schema crudo del
        proveedor a partir de Terraform. El segundo descompone ese schema, que es
        monolítico y de gran tamaño, en una colección de ficheros individuales, uno
        por cada recurso o fuente de datos. El tercero produce a partir de cada schema individual una plantilla de código HCL
        pre-rellena con los atributos del recurso,
        sus tipos, sus restricciones y sus valores por defecto. El cuarto paso realiza una verificación de la cobertura de tipos y la consistencia de la transformación entre el esquema JSON y la plantilla HCL.
        El objetivo es garantizar que no se pierde información durante la traducción de JSON a HCL y que todas las estructuras del esquema original están correctamente reflejadas en el artefacto final.

        \justify
        Cada uno de estos pasos está implementado como un programa independiente
        escrito en Go. La elección de Go se debe a dos
        razones concretas. Primero, existe una biblioteca oficial mantenida por
        HashiCorp para interpretar el formato del schema.
        Segundo, Go es particularmente adecuado para herramientas de procesamiento masivo de ficheros JSON, que es exactamente la
        naturaleza del problema.

        \subsection{Paso 1: obtención del schema crudo}
        \justify
        El punto de partida es la propia herramienta Terraform. Terraform permite invocar un comando que
        serializa todo lo que ese binario sabe sobre sí mismo en un único fichero
        JSON. El proceso no requiere conexión con el proveedor cloud ni credenciales
        de ningún tipo: el schema describe \textit{qué recursos y que atributos tienen}.

        \justify
        El fichero resultante presenta una estructura jerárquica con tres niveles. En
        el nivel raíz figura el conjunto de proveedores cargados, identificados por su
        nombre canónico. Dentro de cada proveedor aparecen cinco categorías de
        entidades: la configuración del propio proveedor, los \textbf{recursos}
        (entidades que el usuario crea), las \textbf{fuentes de datos} (entidades que
        el usuario consulta sin crearlas), los recursos efímeros y las funciones.
        Dentro de cada entidad se encuentra la descripción detallada en forma de
        bloques recursivos: cada bloque puede contener atributos primitivos y
        sub-bloques, que a su vez contienen sus propios atributos y sub-bloques.
        \justify
        Para AWS en su versión 5.x, este paso produce un fichero monolítico de
        aproximadamente \textbf{trece megabytes} que contiene la descripción completa
        del proveedor.

        \subsection{Paso 2: descomposición en schemas individuales}
        \justify
        Trabajar directamente sobre el fichero monolítico no es una opción viable. La solución adoptada es descomponer el fichero monolítico en una colección de
        ficheros individuales, uno por cada recurso o fuente de datos del proveedor,
        organizados en una jerarquía de directorios que refleja la categorización
        original. La aplicación consulta entonces el fichero correspondiente al
        recurso que necesita, sin recorrer el resto. El coste de la descomposición se
        amortiza una sola vez por versión del proveedor, mientras que el beneficio se
        materializa en cada consulta posterior.

        \justify
        La decisión técnica más relevante de este paso es que los ficheros se generan
        \textbf{preservando byte a byte} la información original, sin reinterpretarla
        ni reformatearla semánticamente. Si el proveedor introdujera en una versión
        futura un campo nuevo en la descripción de los atributos, ese campo quedaría
        preservado automáticamente, sin que el extractor tenga siquiera que conocer su
        existencia.

        \justify
        El resultado de este paso, ejecutado sobre el provider de AWS, es una
        jerarquía de \textbf{2.144 ficheros JSON} organizados por categoría: 1.526
        recursos, 608 fuentes de datos, 7 recursos efímeros, 3 funciones y la
        configuración del propio proveedor. Cada fichero conserva exactamente lo que
        el binario exportó originalmente.

        \justify
        Conviene distinguir aquí entre dos conceptos diferentes que conviven en el
        proyecto: la \textbf{capacidad de la pipeline} y el \textbf{catálogo del
        MVP}. La pipeline es capaz de extraer el catálogo completo del proveedor
        las 2.144 entidades, y de hacerlo de forma
        reproducible para cualquier versión futura. El catálogo que la aplicación
        distribuye y muestra al usuario, en cambio, se ha acotado deliberadamente a
        un \textbf{subconjunto representativo} de recursos AWS, seleccionado por cubrir los escenarios habituales de aprendizaje y
        prototipado descritos en el alcance del MVP. Esta separación permite
        validar el flujo completo de la herramienta sobre un catálogo manejable,
        sin que ello impida ampliar la cobertura simplemente añadiendo más
        artefactos extraídos por la pipeline, sin tocar el núcleo de la
        aplicación.

        \subsection{Paso 3: generación de plantillas HCL}
        \justify
        El tercer paso es el más rico desde el punto de vista de diseño. Su objetivo
        es dado el schema JSON de un recurso, producir un fichero
        HCL pre-rellenado que liste todos sus atributos y bloques anidados, con
        valores por defecto sintácticamente válidos y comentarios que documenten
        tipo, restricciones semánticas y cardinalidad.

        \justify
        La utilidad de estas plantillas en el sistema es doble. Por un lado, son la
        base sobre la que la aplicación construye el código IaC real cuando el
        usuario añade un recurso a su diseño: la plantilla aporta el esqueleto
        sintáctico y la aplicación rellena los valores reales que el usuario introduce
        a través de la interfaz visual. Por otro lado, las plantillas son
        consultables como \textbf{material de referencia}: un usuario que se encuentra
        con un recurso desconocido puede inspeccionar su plantilla para entender de un
        vistazo qué atributos admite y con qué tipos, sin necesidad de consultar
        documentación externa.

        \justify
        Para realizar esta traducción es imprescindible interpretar el sistema de
        tipos que utiliza Terraform internamente. A diferencia de un sistema plano,
        este sistema admite tipos primitivos (cadenas, números, booleanos), tipos
        colección (listas, conjuntos y mapas), tipos estructurados (objetos con
        campos nombrados y tuplas con posiciones) y combinaciones recursivas
        arbitrarias entre todos los anteriores. 

        \justify
        Cada atributo de la plantilla se acompaña, además del comentario de tipo, de
        un conjunto de \textbf{flags semánticos} que indican cómo debe tratarlo el
        usuario. Campos marcados como \textit{required} es un campo que el usuario debe proporcionar
        obligatoriamente; aparece en la plantilla con un valor por defecto vacío que
        el usuario debe sustituir. Un atributo \textit{optional} es un campo que el
        usuario puede proporcionar o no; aparece de forma similar pero con la
        anotación correspondiente. La combinación \textit{optional + computed} indica
        que el usuario puede proporcionar el valor o, si no lo hace, el proveedor lo
        calculará automáticamente. La plantilla los presenta comentados, de modo que el
        lector ve que existen y conoce su tipo, pero no afecta a la plantilla.

        \justify
        El generador también traduce a la plantilla la información de
        \textbf{cardinalidad} de los bloques anidados. Cuando el schema indica que un
        bloque debe aparecer un número mínimo o máximo de veces, ese requisito queda
        reflejado en un comentario inline junto a la cabecera del bloque, con
        notaciones del estilo \texttt{[0..1]}, \texttt{[1..*]} o \texttt{[2..5]}.


        \subsection{Paso 4: auditoría de cobertura}
        \justify
        El cuarto y último paso de la pipeline no produce artefactos consumibles, sino
        que verifica la corrección del trabajo realizado por los anteriores. Si en una
        actualización futura del proveedor apareciera un tipo nuevo no contemplado, el
        generador podría procesarlo de forma genérica.

        \justify
        El auditor recorre el conjunto completo de schemas extraídos y clasifica cada
        tipo encontrado contra la lista de tipos conocidos. Cualquier tipo que no
        encaje queda reportado con la ruta exacta donde apareció (recurso y
        atributo), de modo que la adaptación de la pipeline a la nueva versión del
        proveedor pueda hacerse de forma atómica.

        % poner esquema de la pipeline funcionando

        \subsection{Decisiones de diseño relevantes}
        \justify
        El diseño de la pipeline tal como se ha descrito es el resultado de varias
        iteraciones, y algunas decisiones merecen un comentario explícito porque sus
        razones no son obvias y han condicionado el resto del proyecto.

        \justify
        La primera decisión fue \textbf{separar la pipeline en cuatro programas
        independientes} en lugar de unificarlo en uno solo. Aunque la opción
        unificada habría sido más simple a primera vista, mantener los pasos
        separados aporta tres beneficios concretos: cada programa tiene
        \textit{dependencias mínimas} adaptadas a su tarea, los pasos pueden
        \textit{ejecutarse de forma independiente} con frecuencias distintas (la
        extracción solo ante cambios de versión del proveedor; la generación de
        plantillas y la auditoría cuando convenga), y un fallo en uno de los pasos no
        compromete a los demás, lo que mejora la robustez global del sistema.

        \justify
        La segunda decisión fue \textbf{ejecutar la pipeline una sola vez por versión
        del proveedor} y no integrarlo en el flujo de ejecución habitual de la
        aplicación. La generación del catálogo es una operación de \textit{preproceso
        offline} que produce ficheros estáticos, los cuales pasan a formar parte de
        la propia distribución de la aplicación. Esto significa que el usuario final
        no tiene que ejecutar Go ni la pipeline para utilizar la herramienta.

        \subsection{Síntesis}
        \justify
        La pipeline descrita en esta sección resuelve, de forma compacta, cuatro
        problemas a la vez. Proporciona una representación única, fiable y completa
        del proveedor cloud sin depender de su documentación pública. Indexa esa
        representación por nombre de recurso, lo que permite a la aplicación consultar
        cualquier entidad en tiempo prácticamente constante. Genera material legible
        para humanos, en forma de plantillas HCL anotadas, que sirve como referencia
        de cara al usuario. Y verifica, mediante una auditoría automática, que el
        conjunto de tipos está completamente cubierto, alertando si una
        evolución futura del proveedor introdujera elementos no contemplados.

        \justify
        El resultado de este pipeline es un \textbf{catálogo de recursos} que la
        aplicación incorpora como parte de su distribución y que constituye la base de
        todo lo que se describe en las secciones siguientes. Si en
        el futuro se quisiera incorporar soporte para Google Cloud o Microsoft Azure,
        bastaría con ejecutar el mismo pipeline contra los proveedores
        correspondientes y añadir los catálogos resultantes a la distribución, sin
        necesidad de modificar la aplicación. Esta capacidad de extensión sin
        cambios en el núcleo es la materialización concreta del subobjetivo de
        \textit{arquitectura modular y extensible} planteado en la \ref{sec:Alcance_subobjetivos}.
        \par
    }


    \newpage
    \section{El modelo de dominio}
    {
        \setstretch{1.2}
        \setlength{\parskip}{1em}
        \label{sec:modelo_dominio}
        \justify
        Esta sección se organiza en dos partes. La primera presenta el
        \textbf{modelo conceptual}: las entidades canónicas, sus relaciones, las
        reglas de consistencia y el ciclo de vida, formulados con independencia
        de cualquier tecnología y resumidos en el diagrama UML de la
        Figura~\ref{fig:uml_dominio}. La segunda parte describe la
        \textbf{implementación} real de ese modelo en el código y se apoya en el
        diagrama UML de la Figura~\ref{fig:uml_dominio_impl}, donde se aprecian
        las diferencias entre el diseño ideal y su concreción en código.

        \subsection{Funcionamiento del modelo de dominio}
        \justify
        El modelo de dominio representa el estado real del sistema y actúa como
        \textbf{fuente de verdad}. Su objetivo es capturar la semántica de la
        infraestructura diseñada por el usuario de forma \textbf{neutral} respecto a
        la interfaz visual y al formato IaC final. Esta separación permite que la UI
        evolucione sin alterar la lógica y que la generación de IaC cambie sin
        reescribir el núcleo.

        \justify
        A nivel de diseño, el modelo se construye para cumplir cuatro propiedades:
        \textbf{coherencia} (un único estado válido en cada instante), \textbf{trazabilidad}
        (cada cambio puede seguirse desde la acción del usuario hasta su persistencia),
        \textbf{mantenibilidad} (reglas de negocio concentradas y desacopladas de la UI) y
        \textbf{extensibilidad} (posibilidad de incorporar nuevos recursos o proveedores sin
        rediseñar el núcleo).

        \justify
        A partir de estas propiedades se define un núcleo estable que abstrae los elementos
        mínimos del dominio y permite razonar sobre el sistema de forma consistente.

        \subsection{Núcleo canónico del dominio}
        \justify
        La abstracción se organiza en cinco conceptos estables que actúan como vocabulario
        común para todo el sistema:

        \begin{itemize}
        \item \textbf{Proyecto:} contenedor global de trabajo y configuración.
        \item \textbf{Vista:} escenario aislado dentro del proyecto (por ejemplo, un entorno).
        \item \textbf{Recurso:} entidad de infraestructura con identidad estable y tipo lógico.
        \item \textbf{Atributo:} propiedad tipada de un recurso.
        \item \textbf{Relación (enlace):} vínculo semántico entre atributos de recursos.
        \end{itemize}

        \justify
        Este patrón (Proyecto $\rightarrow$ Vista $\rightarrow$ Recurso
        $\rightarrow$ Atributo, con Enlaces entre atributos de recursos) es el
        que utiliza el modelo para describir cualquier arquitectura,
        independientemente de su complejidad o del proveedor cloud al que se
        destine. La Figura~\ref{fig:uml_dominio} representa esta estructura
        como un \textbf{diagrama UML de clases conceptual}, donde se hacen
        explícitas las \textbf{multiplicidades} entre entidades. Más adelante,
        en la Sección~\ref{sec:modelo_dominio_impl}, se compara este modelo
        conceptual con su traducción al código.

        \begin{figure}[H]
            \centering
            \includegraphics[
                width=0.50\textwidth,
                height=14cm,
            ]{umlconcepto.png}
            \caption{Diagrama UML de clases del modelo de dominio (vista conceptual)}
            \footnotesize Fuente: elaboración propia
            \label{fig:uml_dominio}
        \end{figure}

        \justify
        Cada recurso del diseño se representa como
        una instancia de la entidad \textit{Recurso}, que captura la
        información mínima necesaria para reconstruir el bloque
        equivalente en el código IaC y para identificar el recurso
        unívocamente dentro del proyecto.

        \justify
        En esa estructura, los \textbf{atributos} son la representación
        estructurada de los valores que el usuario ha asignado a las propiedades
        del recurso. Internamente se almacenan en un mapa plano de claves a
        valores, donde las claves pueden utilizar un convenio de
        \textit{rutas compuestas} para representar bloques anidados sin recurrir
        a estructuras heterogéneas.

        \justify
        La noción de vista introduce un aislamiento operativo. Esta estructura por vistas
        tiene implicaciones directas en la forma en que el
        sistema persiste el trabajo del usuario. Una propiedad importante del esquema
        de persistencia descrito es que cada vista dispone de un directorio de trabajo
        completamente independiente. Este aislamiento es lo que permite que un mismo
        proyecto pueda contener simultáneamente una vista de desarrollo desplegada
        sobre una cuenta cloud y una vista de producción desplegada
        sobre otra distinta, sin riesgo de mezclas. Es también lo que
        permite que el usuario abra varias vistas a la vez.

        \justify
        En conjunto, estos conceptos constituyen el metamodelo interno. No dependen del proveedor
        cloud ni del lenguaje IaC, y son suficientes para describir una arquitectura
        cloud de manera consistente.

        \subsubsection{Relación entre dominio, presentación y capa de aplicación}
        \justify
        Definido el metamodelo, se concreta cómo interactúan las capas. \textbf{La capa de presentación}
        no modifica directamente el modelo. Su función es capturar eventos de usuario y
        enviarlos como comandos. La \textbf{capa de aplicación} recibe esos comandos y
        ejecuta la lógica de dominio correspondiente.

        \justify
        El flujo general es el siguiente:
        \begin{enumerate}
            \item La interfaz registra una acción del usuario (crear, editar, conectar o
            eliminar).
            \item La capa de aplicación traduce esa acción a una operación de dominio.
            \item El modelo valida invariantes y aplica el cambio si es consistente.
            \item El estado actualizado se propaga a la interfaz para su renderizado.
            \item El nuevo estado se persiste y se usa para regenerar IaC cuando procede.
        \end{enumerate}

        \justify
        Esta arquitectura desacoplada evita que la lógica de negocio quede distribuida entre
        componentes visuales y facilita la evolución independiente de cada capa.


        \subsubsection{Reglas de consistencia e invariantes}
        \justify
        Para mantener ese metamodelo coherente en el tiempo, el modelo impone reglas de
        consistencia que se validan en cada operación:

        \begin{itemize}
            \item \textbf{Unicidad de identidad:} cada recurso mantiene un identificador único
            y estable durante toda su vida.
            \item \textbf{Unicidad nominal por vista:} no se permiten colisiones de nombres
            lógicos entre recursos del mismo tipo dentro de la misma vista.
            \item \textbf{Integridad referencial:} un enlace solo es válido si origen y
            destino existen y sus atributos son compatibles.
            \item \textbf{No ambigüedad estructural:} un recurso no puede quedar simultáneamente
            en estados jerárquicos incompatibles.
            \item \textbf{Persistencia coherente:} toda mutación válida del dominio debe poder
            serializarse sin pérdida semántica.
        \end{itemize}

        \subsubsection{Ciclo de vida de los elementos del dominio}
        \justify
        Estas reglas se aplican durante todo el ciclo de vida de cada entidad, que se
        modela de forma explícita:

        \begin{enumerate}
            \item \textbf{Creación:} al añadir un recurso, se genera su identidad, se inicializan
            atributos y se registra su pertenencia a la vista activa.
            \item \textbf{Actualización:} cambios de atributos, conexiones o pertenencias
            modifican el estado conservando invariantes.
            \item \textbf{Sincronización:} tras cada cambio, se actualizan las proyecciones
            visuales y la representación textual.
            \item \textbf{Eliminación:} al borrar un recurso, se eliminan también referencias
            y enlaces dependientes para evitar estado huérfano.
        \end{enumerate}

        \justify
        Este ciclo de vida asegura que el sistema no entre en estados parciales o
        contradictorios tras operaciones encadenadas.

        \subsection{Implementación del modelo}
        \label{sec:modelo_dominio_impl}
        \justify
        Las subsecciones anteriores describen el \textbf{modelo conceptual}: las cinco
        entidades canónicas (Proyecto, Vista, Recurso, Atributo y Enlace), sus
        invariantes y su ciclo de vida, definidas con independencia de cualquier
        tecnología y resumidas en el diagrama de la Figura~\ref{fig:uml_dominio}.
        En este apartado se describe cómo se traducen esos conceptos al código
        del sistema; el contraste con el diagrama de implementación
        (Figura~\ref{fig:uml_dominio_impl}) permite identificar las diferencias
        entre el diseño conceptual y su realización en código.

        \justify
        El modelo se ha implementado en \textbf{TypeScript} dentro de
        la capa de aplicación. La capa de infraestructura (Rust/Tauri) no replica el
        modelo: actúa exclusivamente como puente hacia el sistema de ficheros y la
        \textit{Terraform CLI}, y recibe el estado ya serializado. De este modo,
        toda la lógica de dominio queda concentrada en un único lugar y se evita
        mantener dos representaciones sincronizadas a ambos lados del puente Tauri.

        \justify
        En el código, las cinco entidades conceptuales se traducen en los tipos que
        recoge la Figura~\ref{fig:uml_dominio_impl}. El \textbf{Proyecto}
        (\texttt{Project}) actúa como contenedor de versión, metadatos,
        configuración y la lista de vistas. La \textbf{Vista}
        (\texttt{ViewSnapshot}) agrupa los recursos, los nodos visuales y las
        aristas pertenecientes a un escenario. El \textbf{Recurso}
        (\texttt{TerraformResource}) encapsula tanto la información semántica (tipo
        lógico, identidad, configuración) como la representación visual asociada en
        el lienzo. El \textbf{Enlace} (\texttt{SerializedEdge}) modela la
        conexión visual entre dos nodos del lienzo y, en su interior, contiene una
        lista de \textit{mappings} de atributos (\texttt{CanvasEdgeMapping}) que
        son los que realmente realizan el vínculo semántico entre propiedades de
        recursos. Por su parte, el \textbf{Atributo} no aparece como una entidad
        independiente con su propia clase: queda representado como entradas de un
        diccionario plano dentro de la configuración del recurso
        (\texttt{ResourceConfig.attributes}), donde la clave admite el convenio de
        rutas compuestas comentado anteriormente.

        \justify
        Comparando ambos diagramas (Figuras~\ref{fig:uml_dominio} y
        \ref{fig:uml_dominio_impl}) se observan dos diferencias relevantes
        respecto al modelo conceptual. En primer lugar, la separación
        dominio/presentación no se traduce en clases físicamente disjuntas: la
        información visual (posición, icono, geometría de las aristas) se aloja
        junto a la información semántica dentro de los mismos tipos. La
        separación se mantiene a nivel \textbf{lógico}: las reglas de dominio
        operan únicamente sobre los campos semánticos y la capa de presentación
        es la única que lee y escribe los campos visuales. En segundo lugar, el
        Enlace tiene una estructura de \textbf{dos niveles}: un enlace visual
        (\texttt{SerializedEdge}) entre dos nodos puede contener varios
        \textit{mappings} de atributos, lo que permite expresar con una única
        arista visual múltiples referencias HCL entre los recursos conectados.

        \begin{figure}[H]
            \centering
            \includegraphics[
                width=0.60\textwidth,
                height=18cm,
            ]{umlintegracion.png}
            \caption{Diagrama UML de clases del modelo de dominio (vista de implementación)}
            \footnotesize Fuente: elaboración propia
            \label{fig:uml_dominio_impl}
        \end{figure}


        \subsection{Persistencia y formato de proyecto}
        \justify
        Una vez definido el comportamiento del modelo, se detalla cómo se almacena y
        recupera el estado del proyecto.
        La persistencia del modelo de datos en el sistema de ficheros del
        usuario combina dos formatos complementarios. Por un lado, el
        modelo completo del proyecto se serializa en un \textbf{fichero
        descriptor con extensión propia}, que contiene en formato JSON
        toda la información estructurada: metadatos, configuración,
        vistas, recursos, nodos y aristas. Este fichero es la unidad que
        el usuario abre, guarda y comparte. Su formato JSON es legible,
        diff-able con herramientas estándar y susceptible de versionarse
        con Git como cualquier otro artefacto de un proyecto de
        software, lo que cumple el requisito de \textbf{portabilidad}.

        \justify
        Mantener esta información dentro del modelo de la vista, en
        lugar de derivarla cada vez del modelo de dominio o del código
        IaC, es lo que permite que la experiencia del usuario sea
        \textit{idempotente}.

        \justify
        Por otro lado, junto al fichero descriptor se mantiene una
        \textbf{carpeta hermana} cuya estructura interna espeja la
        organización lógica del proyecto: dentro de ella, cada vista
        dispone de su propio subdirectorio, que contiene el fichero de
        configuración IaC generado a partir del modelo de la vista y
        cualquier fichero auxiliar que el usuario haya añadido (como
        variables, locales o módulos). Esta carpeta hermana es
        precisamente el directorio de trabajo sobre el que opera la
        herramienta IaC.

        \justify
        Desde el punto de vista de trazabilidad, cada guardado representa un \textit{snapshot}
        consistente del proyecto. Esto facilita auditoría de cambios, integración con control
        de versiones y recuperación de sesiones de trabajo sin pérdida de contexto.

        \subsection{Síntesis}
        \justify
        El conjunto compone un modelo
        explícito, serializable y portable que constituye la fuente
        de verdad del trabajo del usuario, sobre la que se proyectan
        todas las representaciones que la aplicación ofrece.

        \justify
        En síntesis, el modelo de dominio no solo organiza datos: define el contrato lógico
        del sistema. La presentación, la generación de IaC y la integración con herramientas
        externas dependen de él, pero no lo sustituyen. Esta decisión arquitectónica permite
        que el proyecto evolucione de forma controlada y verificable, manteniendo la
        coherencia interna del sistema en todo momento.
    }

    
    \newpage
    \section{Capa de presentación}
    {
        \setstretch{1.2}
        \setlength{\parskip}{1em}
        \label{sec:capa_presentacion}
        \justify
        Tras describir el modelo de dominio en la Sección~\ref{sec:modelo_dominio}
        y el catálogo de recursos que lo alimenta en la
        Sección~\ref{sec:pipeline}, esta sección describe la \textbf{capa de
        presentación}, situada en la parte superior del diagrama de contenedores
        de la Figura~\ref{fig:arquitectura_c4_contenedores}. Su papel dentro de
        la arquitectura es doble: por un lado, actúa como \textbf{proyección
        visual} del modelo de dominio, traduciéndolo a los elementos gráficos
        que el usuario percibe en la interfaz; por otro, funciona como
        \textbf{punto de entrada de las acciones del usuario}, capturando sus
        eventos y reenviándolos como comandos hacia la capa de aplicación, que
        es quien modifica el modelo. En coherencia con la separación de
        responsabilidades descrita en la Sección~\ref{sec:modelo_dominio_impl},
        en este nivel \textbf{no se implementa lógica de negocio}: la
        presentación no decide la validez de una operación, no genera código
        IaC y no se comunica con Terraform.

        \subsection{Estructura general de la interfaz}
        \justify
        La interfaz se organiza en torno a una barra superior y cuatro paneles
        principales que conforman el espacio de trabajo del usuario.

        \begin{itemize}
            \item \textbf{Barra superior:} permite cambiar entre vistas del
            proyecto y aloja los accesos a las acciones principales. En ella se
            sitúan los botones de \texttt{plan}, \texttt{apply} y
            \texttt{destroy}, que actúan como accesos directos a las operaciones
            disponibles sobre la vista activa.
            \item \textbf{Panel izquierdo:} catálogo de recursos arrastrables,
            organizado por proveedor y preparado para búsqueda y filtrado.
            \item \textbf{Lienzo central:} zona principal de edición visual,
            donde se construye la arquitectura mediante arrastrar y soltar.
            \item \textbf{Panel derecho:} inspector de propiedades del recurso
            seleccionado, con los campos editables del modelo.
            \item \textbf{Panel inferior:} zona reservada a la terminal embebida,
            al mapa de objetos y a los logs y mensajes de validación.
        \end{itemize}

        \justify
        Esta disposición responde a un patrón habitual en herramientas de diseño
        visual: el usuario dispone de una zona de catálogo, una zona de edición, una
        zona de propiedades y una zona de retroalimentación. Además, la distribución
        de paneles permite adaptar la interfaz a diferentes flujos de trabajo sin
        alterar la lógica subyacente.


        \begin{figure}[H]
            \centering
            \includegraphics[
                width=0.9\textwidth,
                height=7cm,
                keepaspectratio,
            ]{fotoallUI.png}
            \caption{Vista general de la interfaz con los cinco paneles anotados}
            \footnotesize Fuente: elaboración propia
            \label{fig:ui_overview}
        \end{figure}

        \subsection{Catálogo de recursos}
        \justify
        El panel izquierdo presenta el catálogo de recursos extraído por la pipeline
        descrito en la Sección~\ref{sec:pipeline}. Cada recurso se muestra como un
        elemento arrastrable con información suficiente para reconocerlo de forma
        rápida: icono, nombre legible y nombre técnico de Terraform.

        \justify
        La interfaz incorpora mecanismos de búsqueda y filtrado para reducir la
        complejidad de un catálogo potencialmente amplio. El objetivo es que el
        usuario encuentre con facilidad el recurso deseado sin necesidad de navegar
        manualmente por toda la lista. Desde el punto de vista de la experiencia de
        uso, este panel actúa como punto de entrada al diseño de la arquitectura.


        \begin{figure}[H]
            \centering
            \includegraphics[
                width=0.60\textwidth,
                height=11cm,
                keepaspectratio,
            ]{letpanel.png}
            \caption{panel izquierdo: catálogo de recursos con búsqueda activa}
            \footnotesize Fuente: elaboración propia
            \label{fig:ui_catalogo}
        \end{figure}

        \subsection{Inspector de propiedades}
        \justify
        Cuando un recurso está seleccionado en el lienzo, el panel derecho muestra un
        \textbf{inspector} que lista todos los atributos del recurso con sus valores
        actuales. El inspector es la vía por la cual el usuario edita las propiedades
        sin tocar código.

        \justify
        Para cada atributo, el inspector elige automáticamente el editor más adecuado
        según su tipo:
        \begin{itemize}
        \item \textbf{String/number/bool:} input, textarea, toggle.
        \item \textbf{list(string):} lista de inputs con add/remove.
        \item \textbf{object/set(object(...)):} bloque expandible con sub-inputs
        por campo.
        \item \textbf{computed (solo lectura):} input deshabilitado, con tooltip
        ``Terraform calcula este valor''.
        \end{itemize}


        \begin{figure}[H]
            \centering
            \includegraphics[
                width=1.1\textwidth,
                height=14cm,
                keepaspectratio,
            ]{rightpanel2view.png}
            \caption{panel derecho: inspector mostrando distintos tipos de editor}
            \footnotesize Fuente: elaboración propia
            \label{fig:ui_inspector}
        \end{figure}

        \subsection{Panel inferior: terminal y logs}
        \justify
        El panel inferior aloja tres responsabilidades:

        \begin{itemize}
        \item \textbf{Terminal embebida:} Una PTY real que ejecuta comandos en el
        directorio de la vista actual. El usuario puede invocar \texttt{terraform
        plan}, \texttt{terraform apply}, u otros comandos, y ver la salida en tiempo real.
        La terminal puede además \textbf{desacoplarse en una ventana independiente},
        de modo que el usuario pueda seguir trabajando sobre el lienzo principal
        mientras observa la ejecución de Terraform en paralelo, sin perder espacio
        del panel inferior.

        \item \textbf{Mapa de objetos:} Un visor que representa, de forma
        jerárquica y estructurada, los recursos actuales de la vista con sus
        atributos y relaciones. Cumple un doble propósito: permitir al usuario
        inspeccionar de un vistazo el estado del modelo sin tener que recorrer
        el lienzo recurso por recurso, y servir como vista de depuración cuando
        algo no cuadra entre el diseño visual y el HCL generado. Se actualiza
        automáticamente con cada cambio del modelo de dominio.

        \item \textbf{Logs y validación:} Eventos de validación del código HCL,
        errores de referencia, warnings. Se alimentan desde la capa de aplicación cada
        vez que el usuario modifica el código.
        \end{itemize}

        \justify
        Estos elementos proporcionan retroalimentación inmediata.

        
        \begin{figure}[H]
            \centering
            \includegraphics[
                width=0.9\textwidth,
                height=9cm,
                keepaspectratio,
            ]{bottompanel.png}
            \caption{panel inferior: terminal embebida, object mapper y logs}
            \footnotesize Fuente: elaboración propia
            \label{fig:ui_bottompanel}
        \end{figure}

        \subsection{Editor de código HCL}
        \justify
        Además del canvas, la interfaz incorpora una vista de código que permite
        visualizar la definición \textbf{HCL} generada a partir del diseño. Esta vista
        resulta útil para usuarios que desean revisar el resultado textual de la
        arquitectura sin abandonar el entorno visual.

        \justify
        El editor de HCL forma parte de la capa de presentación porque solo muestra y
        permite editar el contenido textual asociado a la vista activa. La
        interpretación semántica del código, su validación y su sincronización con el
        modelo de dominio corresponden con la capa de aplicación. Desde la experiencia de
        usuario, esta vista actúa como complemento del lienzo y facilita una transición
        gradual entre modelado visual y configuración textual.

        \begin{figure}[H]
            \centering
            \includegraphics[
                width=0.9\textwidth,
                height=12cm,
                keepaspectratio,
            ]{codesection.png}
            \caption{editor de código HCL con un main.tf de ejemplo}
            \footnotesize Fuente: elaboración propia
            \label{fig:ui_editor_hcl}
        \end{figure}

        \subsection{Vistas Diff y Cloud}
        \justify
        La interfaz incluye dos proyecciones adicionales del estado del sistema:
        \textbf{Diff} y \textbf{Cloud}. Ambas se entienden como distintas
        proyecciones del mismo núcleo de información, no como modelos
        independientes. La capa de presentación se limita a mostrarlas de forma
        clara y consistente para que el usuario pueda analizar el estado de la
        infraestructura desde perspectivas complementarias.

        \subsubsection*{Vista Diff}
        \justify
        La vista Diff (Figura~\ref{fig:ui_diff}) representa los cambios previstos
        a partir de la comparación entre el estado actual del proyecto y el
        resultado de una operación de planificación. Esta visualización permite
        identificar de forma rápida qué recursos se crearán, modificarán o
        eliminarán.

        
        \begin{figure}[H]
            \centering
            \includegraphics[
                width=0.85\textwidth,
                height=12cm,
                keepaspectratio,
            ]{Diff2.png}
            \caption{vista Diff: cambios previstos por el último plan}
            \footnotesize Fuente: elaboración propia
            \label{fig:ui_diff}
        \end{figure}

        \subsubsection*{Vista Cloud}
        \justify
        La vista Cloud (Figura~\ref{fig:ui_cloud}), por su parte, muestra el
        estado real recuperado del proveedor después de consultar la
        infraestructura desplegada. Su función es ofrecer una referencia visual
        del entorno existente y facilitar la comparación entre lo diseñado, lo
        generado y lo realmente desplegado.


        \begin{figure}[H]
            \centering
            \includegraphics[
                width=0.85\textwidth,
                height=12cm,
                keepaspectratio,
            ]{Cloud.png}
            \caption{vista Cloud: estado real recuperado del proveedor}
            \footnotesize Fuente: elaboración propia
            \label{fig:ui_cloud}
        \end{figure}

        \subsection{Relaciones entre recursos: aristas y enlaces}
        \justify
        Los recursos cloud intercambian datos entre sí mediante referencias. En la
        presentación, esto se visualiza como \textbf{aristas} (líneas) que conectan dos
        recursos. Cada arista no es decorativa: es un contenedor de información que
        declara uno o más \textbf{enlaces}.

        \justify
        Un enlace especifica con precisión qué atributo del recurso origen
        alimenta qué atributo del recurso destino. Por ejemplo, una arista entre
        una VPC y una subnet puede contener un enlace que indica que el
        atributo \texttt{id} de la VPC alimenta al atributo \texttt{vpc\_id} de
        la subnet. El usuario crea la arista visualmente y el enlace queda
        registrado en el modelo, traduciéndose en la referencia correspondiente
        dentro del HCL.

        \justify
        Cuando el usuario crea un enlace, ocurre una operación doble:
        \begin{enumerate}
        \item El enlace se registra en la arista, haciéndose visible como una entrada
        en la UI.
        \item El modelo de dominio se actualiza: el atributo correspondiente del
        recurso destino se rellena con la expresión que apunta al recurso origen.
        \end{enumerate}

        \justify
        Si el usuario elimina un enlace, la arista se limpia y el atributo del
        recurso destino se vacía, manteniendo coherencia entre lo visual y lo estructural.


        \begin{figure}[H]
            \centering
            \includegraphics[
                width=0.9\textwidth,
                height=12cm,
                keepaspectratio,
            ]{mapping.png}
            \caption{vista Enlace: arista con un enlace configurado entre dos recursos}
            \footnotesize Fuente: elaboración propia
            \label{fig:ui_enlace}
        \end{figure}

        \subsection{Síntesis}
        \justify
        En conjunto, la capa de presentación actúa como un \textit{espejo} del modelo
        de dominio. No toma decisiones sobre la validez de operaciones, no genera
        código, no habla con Terraform. Su única responsabilidad es:

        \begin{enumerate}
        \item \textbf{Observar} cambios en el modelo de dominio y re-renderizar.
        \item \textbf{Capturar} eventos del usuario y enviarlos a la capa de
        aplicación.
        \item \textbf{Visualizar} la estructura del proyecto de forma intuitiva:
        recursos, contenedores, aristas, atributos.
        \end{enumerate}

        \justify
        Esta separación clara permite que la presentación evolucione (cambiar colores,
        layouts, bibliotecas de UI) sin afectar la lógica del dominio. Inversamente,
        cambios en el modelo de dominio no rompen la presentación siempre que la
        interfaz se mantenga.
        \par

    }

    


    \newpage
    \section{Capa de aplicación}
    {
        \setstretch{1.2}
        \setlength{\parskip}{1em}

        \justify
        La capa de aplicación es el núcleo lógico del sistema. Su función es
        implementar la lógica de negocio, procesar los comandos recibidos desde la
        presentación, actualizar el modelo de dominio y coordinar la generación de
        código IaC. Es el puente entre lo visual y lo textual, entre la intención del
        usuario y la realidad de la infraestructura.

        \subsection{Responsabilidades principales}
        \justify
        La capa de aplicación asume un conjunto de responsabilidades bien definidas:
        
        \begin{itemize}
            \item \textbf{Gestión del modelo de dominio:} crea, actualiza y elimina
             recursos, atributos y relaciones (aristas/enlaces) de forma consistente.
            \item \textbf{Sincronización modelo--HCL:} convierte el modelo en
            configuración Terraform válida y, cuando el usuario edita el HCL, reconcilia
            esos cambios con el modelo sin romper su estructura.
            \item \textbf{Persistencia del proyecto:} serializa el estado completo en el
            fichero descriptor y mantiene la carpeta de trabajo de cada vista con su
            \texttt{main.tf} y ficheros auxiliares.
            \item \textbf{Aplicación de reglas semánticas:} evita incoherencias como
            atributos inválidos, referencias rotas o estados intermedios que el usuario
            no debería observar.
            \item \textbf{Coordinación con Terraform:} prepara la ejecución (sync de
            ficheros, credenciales, contexto de la vista) y procesa resultados (plan,
            validación, estado cloud) para reflejarlos en la interfaz.
            \item \textbf{Historial de operaciones:} registra los cambios
            significativos del modelo para permitir al usuario deshacer y
            rehacer acciones de forma reversible.
        \end{itemize}

        \subsection{Procesamiento de comandos}
        \justify
        Cada acción del usuario se traduce en un comando que la capa de aplicación
        recibe para su procesamiento. El comando encapsula la intención del usuario en
        términos de operaciones sobre el modelo de dominio (p. ej., ``crear un recurso
        de tipo \texttt{aws\_vpc} con nombre \texttt{my\_vpc}'').

        \justify
        El procesamiento de comandos sigue un patrón general:
        \begin{enumerate}
            \item Validación: se comprueba que el comando es semánticamente correcto y no viola invariantes.
            \item Ejecución: se actualiza el modelo de dominio según lo especificado por el comando.
            \item Notificación: se informa a la presentación que el modelo ha cambiado para que re-renderice.
            \item Persistencia: se guarda el nuevo estado del proyecto en disco.
            \item Generación: si procede, se actualiza el código IaC asociado a la vista activa.
        \end{enumerate}

        \justify
        Este flujo asegura que cada acción del usuario se refleja de forma consistente
        en el estado interno, en la interfaz y en el código generado.

        \subsubsection*{Catálogo de comandos representativos}
        \justify
        La Tabla~\ref{tab:comandos_aplicacion} recoge los comandos
        representativos que esta capa procesa. Cada comando se describe por su
        intención y por su efecto sobre el modelo de dominio.

        \begin{table}[H]
        \centering
        \small
        \begin{tabular}{p{0.32\textwidth} p{0.6\textwidth}}
        \toprule
        \textbf{Comando} & \textbf{Efecto sobre el modelo} \\
        \midrule
        \textit{Crear recurso} & Añade una entidad recurso al modelo con
        identidad fresca, posición inicial y atributos por defecto tomados del
        catálogo. \\
        \addlinespace
        \textit{Editar atributo} & Actualiza el valor de un atributo del
        recurso, validando contra el tipo declarado en el catálogo. \\
        \addlinespace
        \textit{Renombrar recurso} & Cambia el nombre lógico, verificando la
        unicidad nominal dentro de la vista activa. \\
        \addlinespace
        \textit{Mover recurso} & Actualiza la posición visual del recurso sin
        alterar su semántica. \\
        \addlinespace
        \textit{Eliminar recurso} & Elimina el recurso y, en cascada, todos los
        enlaces que lo referencian, evitando estado huérfano. \\
        \addlinespace
        \textit{Crear enlace} & Registra una nueva arista, valida que origen y
        destino existen y son compatibles, y actualiza el atributo destino con
        la referencia correspondiente. \\
        \addlinespace
        \textit{Eliminar enlace} & Limpia la arista y revierte el atributo
        destino al valor por defecto. \\
        \addlinespace
        \textit{Crear vista} & Añade una vista al proyecto con su directorio de
        trabajo aislado. \\
        \addlinespace
        \textit{Cambiar vista activa} & Cambia el contexto de la sesión; no
        muta el modelo, pero reproyecta presentación y código sobre la vista
        elegida. \\
        \addlinespace
        \textit{Importar HCL} & Parsea el código y lo reconcilia con el modelo
        (sentido HCL $\rightarrow$ modelo de la sincronización bidireccional). \\
        \bottomrule
        \end{tabular}
        \caption{Comandos representativos procesados por la capa de aplicación}
        \label{tab:comandos_aplicacion}
        \footnotesize Fuente: elaboración propia.
        \end{table}


        \justify
        Si en cualquier punto se viola una invariante (por ejemplo, una
        colisión nominal o un valor incompatible con el tipo), la mutación se
        aborta antes de ejecutarse y se devuelve un error a la presentación,
        dejando el modelo en su estado anterior. Esta atomicidad es la que
        garantiza que el proyecto nunca quede en un estado parcialmente
        modificado.

        \subsection{Sincronización bidireccional entre modelo y código}
        \justify
        Un rasgo distintivo de la aplicación es que el modelo visual y el HCL no son
        dos fuentes de verdad independientes. La capa de aplicación mantiene ambos
        sincronizados mediante un ciclo bidireccional:

        \begin{itemize}
            \item \textbf{Modelo $\rightarrow$ HCL:} cualquier cambio en el modelo (nuevo
            recurso, edición de atributos, enlace entre recursos) dispara la
            regeneración del código HCL para la vista activa. La salida es
            determinista: a mismo modelo, mismo HCL.
            \item \textbf{HCL $\rightarrow$ Modelo:} cuando el usuario edita el
            código HCL desde la vista de código, el texto se parsea y se reconcilia
            con el modelo. Se actualizan valores de atributos, pero se preservan
            elementos que el modelo necesita para seguir siendo válido (por ejemplo,
            geometría del lienzo o metadatos internos).
        \end{itemize}

        \justify
        Este enfoque evita la divergencia entre lo visual y lo textual, y permite que
        el usuario alterne entre ambos sin romper la coherencia del proyecto.

        \subsection{Historial de operaciones}
        \justify
        La capa de aplicación mantiene
        un \textbf{historial en memoria} de las operaciones realizadas sobre el
        modelo durante la sesión activa. Este historial habilita las acciones de
        \textit{deshacer} y \textit{rehacer}, que el usuario puede invocar
        directamente desde la interfaz. Cada entrada del historial corresponde a
        un cambio significativo del modelo  creación o eliminación
        de un recurso, edición de un atributo, alta o baja de un enlace, y se almacena con la información mínima necesaria para
        revertirlo. Las acciones puramente visuales, como el desplazamiento de
        un nodo en el lienzo, no se registran como pasos independientes para
        evitar saturar el historial con cambios cosméticos.
        referencia.

        \subsection{Integración con operaciones Terraform}
        \justify
        Cuando el usuario solicita \textit{plan}, \textit{apply}, \textit{destroy} o
        \textit{show}, la capa de aplicación prepara el entorno: sincroniza los
        ficheros del proyecto, valida que la vista activa está lista y delega la
        ejecución al backend nativo. Después, interpreta la salida para actualizar
        el estado visual y el estado real de la infraestructura. De esta forma, el resultado de Terraform se traduce en feedback
        directo para el usuario, sin que tenga que abandonar la herramienta.

        \justify
        Es importante notar que la capa de aplicación \textbf{no invoca
        Terraform directamente}: únicamente prepara el contexto y delega la
        ejecución en la capa de infraestructura, que se ocupa de la mecánica
        del subproceso, del \textit{streaming} de salida y de la gestión de la
        confirmación interactiva. La capa de aplicación se mantiene así libre
        de detalles del sistema operativo y puede razonarse sobre ella de forma
        aislada, lo que también facilita su verificación.

        \subsection{Traducción visual a IaC}
        \label{sec:traduccion_iac}
        \justify
        La traducción visual a IaC es el núcleo operativo de esta capa. A
        partir del modelo de dominio, se genera un \texttt{main.tf} completo
        que Terraform puede ejecutar directamente. La generación es
        determinista e idempotente: el mismo modelo produce siempre el mismo
        HCL, lo que facilita trazabilidad y control de cambios.

        \subsubsection*{Entrada: modelo de dominio y catálogo}
        \justify
        La traducción toma como entrada el conjunto de recursos del modelo
        (tipo, nombre, atributos y relaciones) y el catálogo de schemas
        generado por la pipeline. El catálogo aporta estructura y tipos para el
        inspector y el autocompletado, mientras que el modelo aporta los
        valores concretos y las referencias entre recursos definidas en el
        lienzo.

        \subsubsection*{Reglas de mapeo principales}
        \justify
        La generación sigue reglas deterministas:
        \begin{itemize}
            \item \textbf{Recurso $\rightarrow$ bloque HCL:} cada recurso se
            convierte en un bloque \texttt{resource} (o \texttt{data}) con su
            \texttt{type} y \texttt{name}.
            \item \textbf{Atributos planos $\rightarrow$ asignaciones:} las
            claves simples se serializan como \texttt{clave = valor}.
            \item \textbf{Atributos con rutas $\rightarrow$ sub-bloques:} las
            claves con notación \texttt{a.b.c} se expanden a bloques anidados
            con la estructura correspondiente.
            \item \textbf{Estructura global:} el fichero resultante incluye un
            bloque \texttt{terraform \{ required\_providers \{ ... \} \}} y el
            bloque \texttt{provider 'x' \{ region = ... \}} antes de los
            recursos.
        \end{itemize}

        \subsubsection*{Enlaces y referencias entre recursos}
        \justify
        Las aristas del lienzo representan referencias entre recursos. En el
        modelo, un enlace se almacena como una expresión de Terraform (por
        ejemplo \texttt{aws\_vpc.main.id}). Durante la traducción, ese valor se
        escribe sin comillas para que Terraform lo interprete como referencia y
        no como literal.

        \subsubsection*{Contenedores y atributos implícitos}
        \justify
        La posición de un recurso dentro de contenedores visuales también se
        traduce a HCL. Un nodo colocado dentro de un contenedor jerárquico (por
        ejemplo, una VPC o una subnet) hereda referencias implícitas: se
        completan atributos como \texttt{vpc\_id} o \texttt{subnet\_id} con la
        dirección del recurso padre. En el caso de contenedores de zona (por
        ejemplo, \textit{availability zone} o \textit{security group}), la
        pertenencia se refleja como listas de referencias o valores de zona,
        respetando la semántica de cada recurso. Esta lógica evita que el
        usuario tenga que cablear manualmente relaciones estructurales y
        garantiza que el HCL sea coherente con la organización visual del
        lienzo.

        \subsubsection*{Bloques repetibles y colecciones}
        \justify
        Cuando un atributo es una colección de objetos, se generan bloques repetidos en lugar de una única
        asignación con lista. Esta decisión se alinea con la sintaxis esperada
        por Terraform para \texttt{list(object(...))} y
        \texttt{set(object(...))}, y permite que cada entrada se exprese como
        un bloque anidado independiente.

        \subsubsection*{Normalización de valores}
        \justify
        Antes de escribir el HCL, se descartan valores vacíos o no
        significativos para
        evitar configuraciones inválidas. Los tipos se respetan para producir
        literales correctos: booleanos y números se escriben sin comillas, las
        cadenas se serializan con comillas, y las referencias se escriben como
        expresiones.

        \subsubsection*{Resultado y ejemplo mínimo}
        \justify
        El resultado es un \texttt{main.tf} autocontenido, legible y coherente
        con el diseño visual. La generación es idempotente y conservadora: solo
        se materializa lo que el modelo define, sin introducir ruido ni valores
        por defecto no solicitados.

        \justify
        Si el usuario coloca un \texttt{aws\_vpc} y un \texttt{aws\_subnet}, y
        crea un enlace desde \texttt{aws\_vpc.main.id} hacia
        \texttt{subnet.vpc\_id}, el HCL generado incluye la referencia de forma
        explícita:

        \vspace{1em}
        \begin{lstlisting}[
            language=terraform,
            basicstyle=\footnotesize\ttfamily,
            frame=single,
            framesep=5pt,
            backgroundcolor=\color{gray!5},
            caption={VPC con subnet pública},
            captionpos=b,
            belowcaptionskip=-0.5em,
            label=lst:vpc_subnet,
        ]
        resource "aws_vpc" "main" {
            cidr_block = "10.0.0.0/16"
        }
        resource "aws_subnet" "public" {
            vpc_id     = aws_vpc.main.id
            cidr_block = "10.0.1.0/24"
        }
        \end{lstlisting}
        \begin{center}
        \footnotesize\textit{Fuente: elaboración propia}
        \end{center}


    }

    \newpage
    \section{Capa de infraestructura} {
        \setstretch{1.2}
        \setlength{\parskip}{1em}

        \justify
        La capa de infraestructura es el nivel más bajo del sistema y concentra todo
        lo que depende del sistema operativo: ejecución de binarios, acceso a disco,
        procesos interactivos y lectura de credenciales locales. Se implementa en el
        backend nativo en Rust (Tauri) y expone comandos al frontend mediante
        \textit{invoke}, de forma que la interfaz no tenga acceso directo a recursos
        sensibles. Esta capa no conoce la lógica de negocio del modelo; su misión es
        ejecutar acciones fiables y transmitir eventos de vuelta a la aplicación.

        \subsection{Rol, límites y superficie de comandos}
        \justify
        La capa se materializa en un conjunto de comandos Tauri que encapsulan todo
        el ciclo de Terraform: \texttt{plan}, \texttt{plan -destroy}, \texttt{apply},
        \texttt{destroy}, \texttt{show -json}, validación \texttt{validate -json} y
        diagnósticos mediante LSP. Además, expone comandos auxiliares para confirmar
        ejecuciones interactivas (envío de \texttt{yes/no}), calcular una firma del
        \texttt{terraform.tfstate} y gestionar una terminal embebida. De este modo,
        la capa de aplicación opera siempre sobre una API controlada y predecible.

        \subsection{Ejecución de Terraform por vista}
        \justify
        Cada vista del proyecto se ejecuta en su propio directorio Terraform. Esto
        implica que cada comando se lanza con \texttt{cwd} apuntando a la carpeta de
        la vista activa, evitando mezclar estados entre vistas y permitiendo trabajar
        con entornos distintos (por ejemplo, \textit{dev} y \textit{prod}) dentro del
        mismo proyecto. Antes de cada comando se asegura la sincronización del
        \texttt{main.tf} y de los ficheros auxiliares en disco.

        \subsection{Procesos interactivos y confirmación de \texttt{apply/destroy}}
        \justify
        Los comandos \texttt{apply} y \texttt{destroy} son interactivos: Terraform
        solicita confirmación. La capa de infraestructura mantiene el proceso abierto
        con \texttt{stdin} conectado y guarda una referencia segura al descriptor para
        poder enviar \texttt{yes} o \texttt{no} cuando el usuario confirma desde la UI.
        Este diseño evita que el usuario escriba manualmente en una consola y mantiene
        el flujo dentro de la aplicación.

        \justify
        Además de confirmar, el usuario puede \textbf{cancelar} en cualquier
        momento un proceso en curso. En ese caso, la capa de infraestructura
        envía al subproceso una \textbf{señal de interrupción} (equivalente a
        \texttt{SIGINT}) y, si tras un breve margen el proceso no ha terminado,
        escala a una señal de terminación (\texttt{SIGTERM}). Este mecanismo
        está abstraído por sistema operativo, de modo que el comportamiento es
        equivalente en Linux, macOS y Windows. La cancelación cierra la
        ejecución de forma controlada y devuelve el control a la aplicación sin
        dejar procesos huérfanos.

        \subsection{Streaming de salida en tiempo real}
        \justify
        La salida estándar y de error se transmite en vivo. Se crean lectores
        concurrentes para \texttt{stdout} y \texttt{stderr}, se envían por un canal y
        se emiten eventos Tauri con cada línea.

        \subsection{Ejemplo: flujo Plan $\rightarrow$ Diff $\rightarrow$ Cloud}
        \justify
        Un flujo típico de infraestructura se ejecuta así:
        \begin{enumerate}
            \item El usuario pulsa \texttt{plan}. La capa de infraestructura ejecuta
            \texttt{terraform plan} en el directorio de la vista activa y transmite
            la salida en streaming.
            \item La capa de aplicación interpreta el plan y proyecta los cambios en
            la vista Diff: nodos marcados como \textit{create}, \textit{change} o
            \textit{destroy}.
            \item Tras aplicar los cambios, se consulta \texttt{terraform show -json}
            y se actualiza la vista Cloud con los recursos reales y sus valores
            resultantes (IDs, ARNs, IPs).
        \end{enumerate}

        \justify
        Este ejemplo resume cómo la capa de infraestructura convierte el resultado
        de Terraform en información accionable para el usuario, manteniendo el
        ciclo visual conectado con el estado real.

        \subsection{Gestión de credenciales y aislamiento}
        \justify
        La capa resuelve credenciales AWS justo antes de cada ejecución, sin
        persistirlas en disco. Puede obtenerlas de tres fuentes: variables de entorno,
        perfiles locales (\texttt{\textasciitilde/.aws/credentials}) o un fichero \texttt{.env}.
        Las credenciales se inyectan como variables de entorno del subproceso
        (\texttt{AWS\_ACCESS\_KEY\_ID}, \texttt{AWS\_SECRET\_ACCESS\_KEY},
        \texttt{AWS\_SESSION\_TOKEN}, \texttt{AWS\_DEFAULT\_REGION}) y se descartan al
        finalizar la ejecución. Esta estrategia reduce la exposición y evita que la
        interfaz pueda leer valores sensibles.

        \subsection{Estado real y detección de cambios}
        \justify
        Para el modo Cloud, la capa ejecuta \texttt{terraform show -json} y devuelve
        un listado normalizado de recursos del \textit{state}. La aplicación lo indexa
        y compara con el modelo para marcar cada nodo como \textit{present} o
        \textit{missing}. En paralelo, se calcula una firma ligera del
        \texttt{terraform.tfstate} (longitud y \textit{mtime}) para detectar cambios
        externos sin reparsear el fichero continuamente.

        \subsection{Terminal embebida y comandos auxiliares}
        \justify
        Además de Terraform, la infraestructura expone una PTY real para ejecutar
        comandos dentro del directorio de la vista. Esto habilita la terminal
        embebida, la ventana desacoplada y tareas de depuración avanzadas sin
        abandonar la aplicación, manteniendo el flujo de trabajo en un único lugar.

        \subsection{Seguridad y portabilidad}
        \justify
        El diseño prioriza la separación de responsabilidades: la interfaz no accede
        directamente al sistema de ficheros ni a los binarios, y todas las operaciones
        críticas se concentran en el backend nativo. La capa abstrae diferencias de
        rutas, permisos y ejecución entre sistemas operativos, manteniendo el mismo
        comportamiento en Linux, macOS y Windows.
    }

    % \newpage
    % \section{Pruebas y resultados}
    % {
    %     \setstretch{1.2}
    %     \setlength{\parskip}{1em}
    %     \label{sec:pruebas}

    %     \justify
    %     Esta sección recoge la validación práctica del sistema. El objetivo es demostrar que la herramienta cumple los
    %     subobjetivos del Alcance (Sección~\ref{sec:Alcance_subobjetivos}) sobre
    %     casos de uso representativos y que el código IaC generado es directamente
    %     ejecutable por Terraform sin intervención manual.

    %     \subsection{Estrategia de validación}
    %     \justify
    %     La validación se articula en tres niveles complementarios:

    %     \begin{enumerate}[leftmargin=*, itemsep=2pt]
    %         \item \textbf{Validación funcional manual} mediante la ejecución
    %         extremo a extremo de los casos de uso descritos en la
    %         Sección~\ref{sec:casos_uso}, observando que cada paso produce el efecto
    %         esperado en el modelo, en el lienzo y en el código generado.
    %         \item \textbf{Validación del HCL generado} mediante los propios
    %         comandos de Terraform: \texttt{terraform validate} confirma que el
    %         \texttt{main.tf} es sintáctica y semánticamente correcto, y
    %         \texttt{terraform plan} confirma que el plan se construye sin errores
    %         sobre las credenciales del proveedor.
    %         \item \textbf{Validación del flujo interactivo}: se comprueba que la
    %         aplicación es capaz de gestionar las operaciones interactivas de
    %         Terraform (\texttt{apply} y \texttt{destroy}) enviando \texttt{yes}
    %         por \texttt{stdin} al subproceso, y que el modo \textit{auto-accept}
    %         completa los flujos sin intervención manual del usuario.
    %     \end{enumerate}

    %     \subsection{Caso de prueba 1: VPC + Subnet + EC2 (AWS)}
    %     \justify
    %     Este caso valida el flujo principal de diseño y traducción a IaC sobre una
    %     arquitectura mínima representativa: una \texttt{aws\_vpc}, una
    %     \texttt{aws\_subnet} contenida en la VPC y una \texttt{aws\_instance}
    %     ubicada en la subnet, con dos enlaces (subnet$\rightarrow$VPC e
    %     instancia$\rightarrow$subnet).

    %     % TODO: figura — captura del canvas con VPC + Subnet + EC2
    %     % Contenido que debe mostrar:
    %     %  - Los tres recursos colocados en el lienzo con las dos aristas de enlace
    %     %    visibles.
    %     %  - Idealmente, la subnet visualmente contenida dentro de la VPC.
    %     \begin{figure}[H]
    %         \centering
    %         % \includegraphics[width=0.85\textwidth]{prueba1-canvas.png}
    %         \fbox{\parbox{0.85\textwidth}{\centering\vspace{2.5cm}\textit{[Placeholder: canvas con VPC + Subnet + EC2 conectados]}\vspace{2.5cm}}}
    %         \caption{Pendiente — \textit{canvas del caso de prueba 1}}
    %         \label{fig:prueba1_canvas}
    %     \end{figure}

    %     % TODO: figura — HCL generado del caso de prueba 1
    %     % Contenido que debe mostrar:
    %     %  - El main.tf generado, con los tres bloques resource y las dos referencias
    %     %    vpc_id = aws_vpc.main.id y subnet_id = aws_subnet.public.id.
    %     \begin{figure}[H]
    %         \centering
    %         % \includegraphics[width=0.85\textwidth]{prueba1-hcl.png}
    %         \fbox{\parbox{0.85\textwidth}{\centering\vspace{2.5cm}\textit{[Placeholder: HCL generado del caso de prueba 1]}\vspace{2.5cm}}}
    %         \caption{Pendiente — \textit{HCL generado para VPC + Subnet + EC2}}
    %         \label{fig:prueba1_hcl}
    %     \end{figure}

    %     % TODO: figura — salida de terraform plan del caso de prueba 1
    %     % Contenido que debe mostrar:
    %     %  - Salida de terraform plan en la terminal embebida con los tres recursos
    %     %    marcados como "+ create" y el resumen final "Plan: 3 to add, 0 to change,
    %     %    0 to destroy".
    %     \begin{figure}[H]
    %         \centering
    %         % \includegraphics[width=0.85\textwidth]{prueba1-plan.png}
    %         \fbox{\parbox{0.85\textwidth}{\centering\vspace{2.5cm}\textit{[Placeholder: salida de terraform plan del caso de prueba 1]}\vspace{2.5cm}}}
    %         \caption{Pendiente — \textit{salida de \texttt{terraform plan} del caso de prueba 1}}
    %         \label{fig:prueba1_plan}
    %     \end{figure}

    %     % TODO: figura — vista Cloud tras apply del caso de prueba 1
    %     % Contenido que debe mostrar:
    %     %  - El lienzo en modo Cloud, con los tres recursos marcados como "present"
    %     %    y mostrando sus identificadores reales (vpc-..., subnet-..., i-...).
    %     \begin{figure}[H]
    %         \centering
    %         % \includegraphics[width=0.85\textwidth]{prueba1-cloud.png}
    %         \fbox{\parbox{0.85\textwidth}{\centering\vspace{2.5cm}\textit{[Placeholder: vista Cloud tras apply]}\vspace{2.5cm}}}
    %         \caption{Pendiente — \textit{vista Cloud tras \texttt{apply} del caso de prueba 1}}
    %         \label{fig:prueba1_cloud}
    %     \end{figure}

    %     \justify
    %     \textbf{Resultado:} el HCL generado pasa \texttt{terraform validate} sin
    %     errores, el plan se ejecuta correctamente sobre las credenciales AWS del
    %     usuario, y la vista Cloud muestra los tres recursos como
    %     \textit{present} con sus identificadores reales. La equivalencia entre lo
    %     diseñado y lo desplegado se verifica visualmente sin pasos manuales fuera
    %     de la aplicación.

    %     \subsection{Caso de prueba 2: edición bidireccional modelo \texorpdfstring{$\leftrightarrow$}{<->} HCL}
    %     \justify
    %     Este caso valida la propiedad clave de \textbf{sincronización bidireccional}
    %     entre el modelo de dominio y el HCL. Se ejecutan dos escenarios encadenados:
    %     primero se modifica un atributo desde el inspector y se observa que el HCL
    %     se regenera; después se modifica el mismo atributo manualmente en el
    %     editor de código y se observa que el modelo, el lienzo y el inspector
    %     reflejan el cambio.

    %     % TODO: figura — antes/después del editor HCL y del canvas
    %     % Contenido que debe mostrar (idealmente dos capturas lado a lado):
    %     %  - "Antes": editor con un valor original (por ejemplo, cidr_block = "10.0.0.0/16")
    %     %    y el inspector mostrando el mismo valor.
    %     %  - "Después": editor con el valor modificado manualmente y el inspector
    %     %    reflejando el nuevo valor, demostrando la sincronización.
    %     \begin{figure}[H]
    %         \centering
    %         % \includegraphics[width=0.95\textwidth]{prueba2-bidireccional.png}
    %         \fbox{\parbox{0.95\textwidth}{\centering\vspace{2.5cm}\textit{[Placeholder: antes/después de la edición bidireccional]}\vspace{2.5cm}}}
    %         \caption{Pendiente — \textit{sincronización bidireccional modelo $\leftrightarrow$ HCL}}
    %         \label{fig:prueba2_bidireccional}
    %     \end{figure}

    %     \justify
    %     \textbf{Resultado:} ambos sentidos de la sincronización se completan sin
    %     pérdida de información. La geometría del lienzo y los metadatos internos del
    %     modelo se preservan durante la reconciliación con el HCL editado, lo que
    %     confirma que el editor de código y el lienzo son dos proyecciones
    %     consistentes del mismo modelo de dominio.

    %     \subsection{Caso de prueba 3: ciclo plan \texorpdfstring{$\rightarrow$}{->} apply \texorpdfstring{$\rightarrow$}{->} destroy}
    %     \justify
    %     Este caso valida el ciclo de vida completo de un despliegue. Sobre la
    %     arquitectura del caso de prueba 1, se ejecuta \texttt{plan}, se aplica el
    %     despliegue, se modifica el modelo (cambio de \texttt{cidr\_block} en la
    %     subnet) para forzar un nuevo \texttt{plan} con cambios, y finalmente se
    %     ejecuta \texttt{destroy}.

    %     % TODO: figura — vista Diff coloreada
    %     % Contenido que debe mostrar:
    %     %  - El lienzo en modo Diff tras el segundo plan, con los recursos coloreados
    %     %    según el tipo de cambio: el recurso modificado en color "change", los
    %     %    no afectados en "no-op", y al menos uno marcado como "destroy" en el
    %     %    plan final de destrucción.
    %     \begin{figure}[H]
    %         \centering
    %         % \includegraphics[width=0.85\textwidth]{prueba3-diff.png}
    %         \fbox{\parbox{0.85\textwidth}{\centering\vspace{2.5cm}\textit{[Placeholder: vista Diff con nodos coloreados]}\vspace{2.5cm}}}
    %         \caption{Pendiente — \textit{vista Diff durante el ciclo plan/apply/destroy}}
    %         \label{fig:prueba3_diff}
    %     \end{figure}

    %     \justify
    %     \textbf{Resultado:} el ciclo se completa sin que el usuario tenga que
    %     abandonar la aplicación. La vista Diff refleja correctamente cada estado
    %     intermedio del plan, y la vista Cloud confirma tanto la creación como la
    %     eliminación final de los recursos en el proveedor.

    %     \subsection{Caso de prueba 4: auto-accept y confirmación interactiva}
    %     \justify
    %     Tanto \texttt{terraform apply} como \texttt{terraform destroy} son comandos
    %     interactivos que solicitan confirmación explícita antes de actuar. Este
    %     caso valida que la aplicación intercepta esa confirmación y permite
    %     responderla desde la UI sin abrir una terminal externa, así como el
    %     funcionamiento del modo \textit{auto-accept}.

    %     % TODO: figura — confirmación interactiva resuelta desde la UI
    %     % Contenido que debe mostrar:
    %     %  - El prompt de Terraform "Enter a value:" detectado por la aplicación,
    %     %    con un diálogo o control en la UI que ofrece los botones "yes" y "no".
    %     %  - Idealmente, un segundo panel mostrando el modo auto-accept activado
    %     %    (un toggle o ajuste) y la operación completándose sin intervención.
    %     \begin{figure}[H]
    %         \centering
    %         % \includegraphics[width=0.85\textwidth]{prueba4-autoaccept.png}
    %         \fbox{\parbox{0.85\textwidth}{\centering\vspace{2.5cm}\textit{[Placeholder: confirmación interactiva y auto-accept]}\vspace{2.5cm}}}
    %         \caption{Pendiente — \textit{confirmación interactiva de Terraform resuelta desde la aplicación}}
    %         \label{fig:prueba4_autoaccept}
    %     \end{figure}

    %     \justify
    %     \textbf{Resultado:} el envío de \texttt{yes} por \texttt{stdin} al
    %     subproceso desbloquea correctamente la ejecución de \texttt{apply} y
    %     \texttt{destroy}. El modo \textit{auto-accept} reduce el ciclo de
    %     despliegue a un único clic, sin comprometer la trazabilidad: la salida del
    %     subproceso sigue mostrándose en tiempo real en la terminal embebida.

    %     \subsection{Discusión}
    %     \justify
    %     En conjunto, los cuatro casos de prueba evidencian el cumplimiento de los
    %     subobjetivos definidos en el Alcance:

    %     \begin{itemize}[leftmargin=*, itemsep=2pt]
    %         \item La \textbf{interfaz visual} permite diseñar arquitecturas no
    %         triviales con varios recursos y enlaces (CP-1, CP-2).
    %         \item La \textbf{traducción automática a IaC} produce HCL ejecutable
    %         que pasa \texttt{terraform validate} y \texttt{terraform plan} (CP-1,
    %         CP-3).
    %         \item La \textbf{arquitectura modular} se materializa en la
    %         integración funcional de los cuatro componentes (pipeline, modelo,
    %         aplicación, infraestructura) durante todo el ciclo (CP-3, CP-4).
    %         \item Los \textbf{casos de uso representativos} validan que el sistema
    %         soporta el flujo completo desde el diseño hasta la destrucción de la
    %         infraestructura desplegada (CP-3, CP-4).
    %     \end{itemize}

    %     \justify
    %     Las limitaciones detectadas durante esta validación son coherentes con el
    %     alcance de un MVP y se recogen como trabajo futuro en la sección de
    %     Conclusiones.
    %     \par
    % }

    \newpage
    \section{Conclusiones}
    {
        \setstretch{1.2}
        \setlength{\parskip}{1em}
        \label{sec:conclusiones}

        \subsection{Conclusiones generales del proyecto}
        \justify
        Este Trabajo Final de Grado planteó el diseño y desarrollo de un sistema
        visual para la creación y despliegue de arquitecturas cloud mediante
        Infraestructura como Código, con un enfoque open source y neutral
        respecto al proveedor. A lo largo del documento se ha cubierto la
        justificación del problema, el alcance, la planificación, el diseño
        arquitectónico, la implementación de las cuatro capas, la pipeline de
        extracción de recursos y la validación funcional sobre casos
        representativos.

        \justify
        El resultado entregado es una aplicación de escritorio funcional capaz de
        modelar visualmente arquitecturas cloud, generar el HCL correspondiente
        de forma determinista y ejecutar el ciclo completo
        \textit{plan} $\rightarrow$ \textit{apply} $\rightarrow$ \textit{destroy}
        sobre un proveedor real, sin que el usuario tenga que abandonar el
        entorno visual. La aplicación incorpora un catálogo extraído de forma
        automática a partir del schema oficial de Terraform, lo que garantiza
        fidelidad y abre la puerta a la incorporación de nuevos proveedores sin
        modificar el núcleo de la aplicación.

        \justify
        El proyecto, planteado como un MVP, alcanza sus objetivos y deja
        sentadas las bases sobre las que pueden
        construirse las evoluciones descritas más adelante.

        \subsection{Evaluación de la propuesta}
        \justify
        Los cuatro subobjetivos definidos en la
        Sección~\ref{sec:Alcance_subobjetivos} se han materializado de forma
        verificable en la solución entregada:

        \begin{itemize}[leftmargin=*, itemsep=2pt]
            \item \textbf{Interfaz gráfica para modelar arquitecturas cloud}:
            se ha logrado una interfaz visual completa, descrita en la
            Sección~\ref{sec:capa_presentacion}, que permite al usuario diseñar
            arquitecturas mediante arrastrar y soltar recursos desde un catálogo,
            editar sus atributos en un inspector contextual y conectarlos entre sí
            mediante enlaces. La interfaz integra además un editor HCL embebido,
            una terminal y vistas Diff y Cloud que proporcionan retroalimentación
            inmediata sobre el estado del proyecto y de la infraestructura
            desplegada.
            \item \textbf{Traducción automática a IaC}: se describe en la sección
            de traducción visual a IaC y se valida con la ejecución de
            \texttt{terraform validate} y \texttt{terraform plan} sobre el HCL
            generado.
            \item \textbf{Arquitectura modular y extensible}: se materializa en
            la pipeline de extracción (Sección~\ref{sec:pipeline}) y en el modelo
            de dominio (Sección~\ref{sec:modelo_dominio}), que aíslan el catálogo
            del proveedor del núcleo de la aplicación y permiten incorporar
            nuevos proveedores sin modificarlo.
        \end{itemize}

        \justify
        Respecto al estado del arte, la propuesta se sitúa en el hueco
        identificado en la
        Tabla~\ref{tab:comparativa_soluciones}: combina modelado visual,
        independencia de proveedor, licenciamiento open source, generación de
        IaC editable y ejecución local. Las soluciones existentes cubren
        subconjuntos parciales de esos criterios, pero ninguna los reúne
        simultáneamente, lo que confirma la pertinencia de la propuesta.

        \subsection{Limitaciones y trabajo futuro}
        \justify
        Aun cumpliendo los objetivos planteados, el alcance del MVP deja
        limitaciones explícitas y abre varias líneas naturales de evolución.

        \subsubsection*{Limitaciones actuales}
        \begin{itemize}[leftmargin=*, itemsep=2pt]
            \item \textbf{Distribución limitada a Linux}: el desarrollo y la
            validación se han realizado sobre Linux. Aunque Tauri soporta
            macOS y Windows de forma nativa, el empaquetado, la firma y la
            distribución estable en esas plataformas requieren trámites
            (certificados, instaladores, actualizaciones) ajenos al núcleo
            técnico del proyecto.
            % \item \textbf{Validación manual mediante casos de uso}: la
            % verificación se ha realizado ejecutando los casos descritos en
            % la Sección~\ref{sec:pruebas} de forma manual. No se ha construido
            % una batería de pruebas automatizadas, coherente con el carácter
            % de prototipo del MVP.
            \item \textbf{Análisis estático de seguridad pendiente}: la
            planificación inicial contemplaba complementar la validación
            funcional con un análisis estático del HCL generado mediante
            herramientas como Trivy, orientado a detectar configuraciones
            inseguras (por ejemplo, grupos de seguridad demasiado abiertos o
            buckets públicos por defecto). Esa integración no se ha podido
            completar dentro del marco temporal del proyecto y queda como
            tarea pendiente, no como línea de evolución del producto.
            \item \textbf{Soporte multi-proveedor.} La pipeline está diseñada
            para ejecutarse sobre cualquier proveedor de Terraform; queda
            pendiente generar los catálogos de Google Cloud y Microsoft Azure
            y adaptar las reglas de mapeo específicas de cada uno.
            \item \textbf{Flujos avanzados de Terraform fuera del MVP}:
            \textit{workspaces}, \textit{backends} remotos y mecanismos
            similares no están integrados en la experiencia visual. Son
            funcionalidades operativas más cercanas a entornos de producción
            que al objetivo de aprendizaje y prototipado del MVP.
        \end{itemize}

        \subsubsection*{Líneas de trabajo futuro}
        \begin{itemize}[leftmargin=*, itemsep=2pt]
            \item \textbf{Integración con un servidor MCP \textit{(Model
            Context Protocol)}}~\cite{noauthor_what_nodate}\textbf{.} Exponer la aplicación como servidor MCP
            permitiría que asistentes de IA pudieran inspeccionar el modelo
            de dominio, proponer cambios sobre la arquitectura visual y
            consumir el catálogo de recursos directamente, abriendo un flujo
            de trabajo asistido sin sacrificar el control del usuario.
            \item \textbf{Sincronización del proyecto con un repositorio
            Git.} El formato de proyecto ya es legible y diff-able; falta
            integración nativa para vincular cada vista con un repositorio,
            registrar commits automáticos en puntos significativos del ciclo
            de vida y permitir colaboración multi-usuario con resolución
            asistida de conflictos.
            \item \textbf{Soporte para módulos Terraform.} Permitiría
            encapsular sub-arquitecturas reutilizables como nodos compuestos
            en el lienzo, aumentando la productividad en proyectos grandes.
            \item \textbf{Importación de infraestructura existente.} Un
            asistente para \texttt{terraform import} permitiría reconstruir
            el modelo a partir de una infraestructura ya desplegada,
            ampliando el público objetivo a equipos con infraestructura
            preexistente.
            \item \textbf{Ejecución automatizada de la pipeline de
            extracción.} Actualmente la pipeline se ejecuta de forma manual
            ante cada cambio de versión del proveedor y los catálogos
            resultantes se importan al repositorio de la aplicación. Como
            evolución natural, se plantea integrarla en un flujo de
            GitHub Actions~\cite{noauthor_descripcion_nodate-1} que detecte nuevas versiones de los proveedores
            de Terraform (o se dispare de forma periódica), ejecute la
            pipeline completa de forma reproducible y abra un \textit{pull
            request} con los catálogos y plantillas actualizadas. Esto
            elimina por completo el paso manual de importación y mantiene
            el catálogo distribuido sincronizado con la última versión
            estable de cada proveedor.
        \end{itemize}

        \subsection{Valoración personal}
        \justify
        En lo personal, el balance del proyecto es muy positivo. Una de las
        motivaciones iniciales era \textbf{profundizar en Terraform y en el
        ecosistema Cloud} desde un ángulo distinto, y
        construir esta herramienta me ha obligado a entender Terraform a un
        nivel que difícilmente habría alcanzado limitándome a usarlo: leer y
        procesar su schema, integrar su CLI como subproceso, gestionar sus
        flujos interactivos y reconstruir el sentido de su sintaxis para
        traducirla desde un modelo visual.

        \justify
        En paralelo, el proyecto me ha permitido \textbf{aprender y consolidar
        Rust} como lenguaje de sistemas para el backend nativo, y \textbf{ganar
        soltura en el desarrollo frontend} moderno con React y TypeScript.
        Llegar a tener las dos partes comunicándose de forma fluida a través de
        Tauri, con una capa de infraestructura sólida en Rust, ha sido uno de
        los aprendizajes más valiosos del TFG.

        \justify
        Más allá del aprendizaje técnico, una motivación importante ha sido
        \textbf{aportar algo a la comunidad open source del Cloud}. Las
        herramientas visuales actuales son mayoritariamente propietarias o
        están ligadas a un proveedor; poder contribuir con una alternativa
        abierta, neutral y orientada al aprendizaje, aunque sea como punto de
        partida, es algo que me satisface en lo profesional y en lo personal.
    }

    \newpage
    \printbibliography

\end{document}
